\documentclass[11pt,a4paper]{article}
\usepackage[T1]{fontenc}
\usepackage[utf8]{inputenc}
\usepackage{lmodern}
\usepackage[a4paper,margin=25mm,headheight=14pt]{geometry}
\usepackage{amsmath,amssymb,amsthm,mathtools,bm}
\usepackage{booktabs,longtable,array}
\usepackage{microtype}
\usepackage{xcolor}
\usepackage{enumitem}
\usepackage{xurl}
\usepackage{fancyhdr}
\usepackage{needspace}
\usepackage{listings}
\usepackage{hyperref}
\hypersetup{colorlinks=true,linkcolor=blue!45!black,citecolor=green!30!black,urlcolor=blue!55!black,
 pdftitle={Combinatorial Transseries and Their Inverses},
 pdfauthor={Research extension prepared for the ProveIt series},
 pdfsubject={Asymptotic expansions and inverse functions of combinatorial sequences}}
\pagestyle{fancy}
\fancyhf{}
\fancyhead[L]{\small Combinatorial transseries and their inverses}
\fancyhead[R]{\small\thepage}
\renewcommand{\headrulewidth}{0.3pt}
\setlength{\parskip}{4pt}
\setlength{\parindent}{1.2em}
\setlength{\emergencystretch}{2em}
\allowdisplaybreaks[2]
\numberwithin{equation}{section}
\newtheorem{theorem}{Theorem}[section]
\newtheorem{proposition}[theorem]{Proposition}
\newtheorem{lemma}[theorem]{Lemma}
\newtheorem{corollary}[theorem]{Corollary}
\theoremstyle{definition}
\newtheorem{definition}[theorem]{Definition}
\newtheorem{example}[theorem]{Example}
\theoremstyle{remark}
\newtheorem{remark}[theorem]{Remark}
\newcolumntype{P}[1]{>{\raggedright\arraybackslash}p{#1}}
\newcommand{\N}{\mathbb N}
\newcommand{\R}{\mathbb R}
\newcommand{\C}{\mathbb C}
\newcommand{\dd}{\,\mathrm d}
\newcommand{\ee}{\mathrm e}
\newcommand{\ii}{\mathrm i}
\newcommand{\coeff}[2]{\left[#1\right]#2}
\newcommand{\cE}{\mathcal E}
\newcommand{\cL}{\mathcal L}
\newcommand{\bigO}{\mathcal O}
\newcommand{\sectorsim}{\mathrel{\sim_{\mathrm{sec}}}}
\newcommand{\fall}[2]{#1^{\underline{#2}}}
\newcommand{\rising}[2]{(#1)_{#2}}
\newcommand{\StirII}[2]{\left\{\begin{matrix}#1\\#2\end{matrix}\right\}}
\newcommand{\StirI}[2]{\left[\begin{matrix}#1\\#2\end{matrix}\right]}
\DeclareMathOperator{\Li}{Li}
\DeclareMathOperator{\ord}{ord}
\DeclareMathOperator{\Rea}{Re}
\DeclareMathOperator{\sech}{sech}
\lstset{basicstyle=\ttfamily\small,breaklines=true,columns=fullflexible,keepspaces=true,
 frame=single,framerule=0.3pt,rulecolor=\color{black!30},showstringspaces=false}

\begin{document}
\begin{titlepage}
\centering
\vspace*{13mm}
{\small A COMPANION TO \textit{TRANSSERIES AND INVERSION}}\par
\vspace{12mm}
{\LARGE\bfseries Combinatorial Transseries\\[3mm]and Their Inverses\par}
\vspace{7mm}
{\large Gamma quotients, multiple saddles,\\arithmetic sectors, and discrete thresholds\par}
\vspace{10mm}
{\normalsize Research extension prepared for the ProveIt series\par}
{\normalsize September 4, 2026\par}
\vspace{12mm}
\begin{minipage}{0.91\textwidth}
\begin{abstract}
The asymptotic inversion of a combinatorial sequence requires more than reversing its leading term. One must identify the correct asymptotic scale, choose or avoid a continuous interpolation, retain subdominant contributions when they are meaningful, and control the amplification of forward errors by inversion. This article develops these points through complementary families not central to the linked \emph{Transseries and Inversion} volume: Catalan and Fuss--Catalan numbers, central multinomial coefficients and rectangular tableaux, both kinds of fixed-column Stirling numbers, Motzkin and Delannoy numbers, large Schr\"oder numbers, involutions, alternating permutations, connected labeled graphs, necklaces, Lyndon words, and harmonic-number functions.

General finite coefficient formulas are given for gamma-quotient expansions, moving-saddle expansions, multiexponential inverses, and inverse sector interactions. Exact moment decompositions expose the alternating subdominant terms of Motzkin numbers and involutions. A fixed-order theorem gives every exponential layer of connected-graph enumeration, including a vanishing forward layer that is regenerated by inversion. Necklace divisor sums demonstrate why arithmetic indicators and accumulating actions cannot be replaced by an ordinary finite exponential grid. Harmonic inversion supplies a contrasting exponential and endpoint-Puiseux regime.

The results are separated into convergent identities, all-orders Poincar\'e expansions, exact sector decompositions with asymptotic amplitudes, and formal inverses valid on a sequence's range. Proofs, explicit coefficients, numerical checks, and reproducible Python code accompany the theory. No global uniqueness of interpolation, blanket convergence of divergent expansions, or uncomputed Stokes constants is assumed.
\end{abstract}
\end{minipage}
\vfill
{\small Self-contained \LaTeX\ source and verification program accompany this PDF.\par}
\end{titlepage}

\tableofcontents
\newpage

\section{Purpose, scope, and the meaning of an inverse}
\label{sec:scope}

\subsection{Relation to the linked article}

The reference volume \emph{Transseries: the polynomial--logarithmic calculus, series reversal at infinity, and the inversion of rapidly growing functions} develops a broad calculus of asymptotic series, Lambert-centered reversal, flat terms, and inverse functions \cite{base}. Its current source also treats such examples as Bell and Fubini numbers, partitions, rooted unlabeled trees, derangements, and several factorial-related functions. Repeating those case studies would add little. Here the emphasis is instead on several additional mechanisms:
\begin{enumerate}[label=(\roman*),itemsep=2pt]
\item balanced gamma quotients with a negative logarithmic exponent, for which the large branch is controlled by $W_{-1}$;
\item finite exponential sums, exact moment decompositions, and two real saddles, whose small sectors can be explicitly normalized;
\item factorial growth with prime-indexed corrections, superexponential graph growth, and divisor-dependent arithmetic corrections;
\item slow growth and finite limiting values, which reverse into exponential and Puiseux scales.
\end{enumerate}
The basic inversion principles are restated because each application needs precise hypotheses. The tables and coefficient formulas are designed to be usable without navigating the much larger parent volume.

A useful concordance is that the parent's source labels \nolinkurl{p0:thm:lambert-core}, \nolinkurl{p0:thm:lambert-centered}, and \nolinkurl{p0:thm:LB} correspond to the Lambert core, centered reversion, and Lagrange--B\"urmann themes used below. Its discussion at \nolinkurl{p0:sec:discrete} is extended by the explicit interpolation counterexamples of Section~\ref{subsec:gauge}. These are source-label references, not claims that the present theorems have been formalized in Lean. The linked editable source and repository description were consulted on September 4, 2026; this companion does not assert page-for-page parity with an older PDF in that directory.

\subsection{Four logically different objects}

Let $(a_n)$ be eventually positive and strictly increasing. The notation ``inverse of $a_n$'' can refer to four different objects.

\begin{definition}
The \emph{inverse on the range} sends $a_n$ to $n$. The \emph{threshold inverse} is the staircase
\begin{equation}
 N(y)=\min\{n\ge n_0:a_n\ge y\}.
 \label{eq:staircase}
\end{equation}
A \emph{specified continuous inverse} is $F^{-1}$ for an explicitly given, eventually increasing function satisfying $F(n)=a_n$. A \emph{formal asymptotic inverse} is a formal expression whose substitution reverses the chosen forward asymptotic model to every prescribed order.
\end{definition}

The four notions often cooperate, but they are not interchangeable. A formal inverse may approximate the range inverse without specifying values of $F$ between integers. An interpolation may be natural from a gamma quotient, a moment integral, or a finite exponential sum. Nevertheless, ``natural'' is a choice, not a uniqueness theorem.

Throughout, $F^{-1}$ means a compositional inverse, never the reciprocal $1/F$. Unless explicitly stated otherwise, $x\to+\infty$, the parameters such as a tableau height or a Stirling column are fixed, and logarithms are real natural logarithms. All assertions of eventual positivity or monotonicity allow discarding a finite initial interval.

\subsection{Interpolation is a mathematical datum}
\label{subsec:gauge}

\begin{proposition}[Interpolation gauge freedom]
\label{prop:gauge}
Let $F$ be an eventually strictly increasing $C^1$ function. There exist other eventually strictly increasing $C^1$ functions with the same values at every sufficiently large integer whose inverse differs from $F^{-1}$ by a nonvanishing bounded oscillation. There also exist such functions whose inverse differs first at any prescribed inverse-power order, or at an exponentially small order while preserving all inverse-power coefficients.
\end{proposition}

\begin{proof}
For $|\eta|<1/(2\pi)$, define
\[
 h(x)=x+\eta\sin(2\pi x),\qquad \widetilde F=F\circ h.
\]
Then $h(n)=n$ and $h'(x)=1+2\pi\eta\cos(2\pi x)>0$. Thus $\widetilde F(n)=F(n)$ and
\[
 \widetilde F^{-1}=h^{-1}\circ F^{-1}.
\]
The map $h^{-1}(u)-u$ is nonzero and periodic if $\eta\ne0$, so the inverse difference need not even tend to zero.

For the finer assertions replace the sine term by $\eta x^{-r}\sin(2\pi x)$, $r>0$, or by $\eta\ee^{-\rho x}\sin(2\pi x)$, $\rho>0$. In both cases $h'(x)>0$ eventually. Direct substitution in $h(h^{-1}(u))=u$ gives, respectively,
\[
 h^{-1}(u)=u-\eta u^{-r}\sin(2\pi u)+o(u^{-r}),
\]
and
\[
 h^{-1}(u)=u-\eta\ee^{-\rho u}\sin(2\pi u)
           +\bigO(\ee^{-2\rho u}).
\]
The latter change is smaller than every inverse power. All the integer samples are nevertheless unchanged.
\end{proof}

Consequently, interpolation-independence of an algebraic inverse expansion holds only within a class of \emph{asymptotically compatible} interpolations. It does not follow from monotonicity and agreement on the integers. This qualification becomes indispensable when exponentially small corrections are discussed.

If $F$ is an exact increasing interpolation, then, on its eventual range,
\begin{equation}
 N(y)=\left\lceil F^{-1}(y)\right\rceil.
 \label{eq:ceil}
\end{equation}
An approximation $\widehat x$ with $|\widehat x-F^{-1}(y)|\le\varepsilon$ determines this ceiling only when the two endpoints $\widehat x\pm\varepsilon$ have the same ceiling. At a sample value $y=a_n$, an inverse approximation with error less than $1/2$ recovers $n$ by nearest-integer rounding. An unqualified ceiling of an $o(1)$ approximation is not a valid threshold algorithm near a jump.

\subsection{What ``transseries'' will mean here}

We distinguish an inverse-power expansion
\[
 F(x)\sim F_0(x)\sum_{j\ge0}a_jx^{-j}
\]
from a sector expansion such as
\[
 F(x)\sectorsim F_0(x)\sum_{k\ge0}\ee^{-k\omega x}
       \sum_{j\ge0}a_{k,j}x^{-j}.
\]
The latter notation means that the separate sector amplitudes have the stated expansions in a specified exact decomposition or a specified formal model. It does \emph{not} mean that truncating the first sector after ten inverse powers leaves an exponentially small error.

A finite additive semigroup of positive actions is locally finite below each action cutoff. This makes coefficient extraction at a fixed exponential order finite. Prime-generated actions remain locally finite because there are finitely many integers below a fixed bound. Necklace actions, in contrast, accumulate at a finite value; Section~\ref{sec:necklaces} explains the resulting limitation.

Factors such as $\cos(\pi x)$ are bounded oscillatory coefficients. They are not positive monomials of an ordinary real logarithmic--exponential Hardy field. Whenever they occur, we work with the explicitly specified real oscillatory extension, or equivalently with conjugate complex exponential sectors. Calling this distinction out avoids assigning an oscillating term a nonexistent eventual sign.

\section{A coefficient and inversion calculus}
\label{sec:calculus}

\subsection{Finite partition formulas}

For a formal series $Q(t)=\sum_{j\ge1}q_jt^j$, set
\begin{equation}
 \cE_n(q)=
 \sum_{\substack{m_1,\ldots,m_n\ge0\\\sum j m_j=n}}
       \prod_{j=1}^n\frac{q_j^{m_j}}{m_j!},
 \qquad \cE_0(q)=1.
 \label{eq:E}
\end{equation}
Then $\exp Q(t)=\sum_{n\ge0}\cE_n(q)t^n$. Thus every coefficient is a finite polynomial, even when the power series itself diverges. For example,
\[
 \cE_1=q_1,\quad
 \cE_2=q_2+\frac{q_1^2}{2},\quad
 \cE_3=q_3+q_1q_2+\frac{q_1^3}{6}.
\]
Conversely, if $A(t)=1+\sum_{j\ge1}a_jt^j$, then
\begin{equation}
 \cL_n(a):=[t^n]\log A(t)
 =\sum_{\substack{\sum j m_j=n\\M=\sum m_j}}
   (-1)^{M+1}(M-1)!\prod_{j=1}^n\frac{a_j^{m_j}}{m_j!}.
 \label{eq:L}
\end{equation}
Both formulas follow by expanding the exponential or logarithm and applying the multinomial theorem. They are the ordinary-coefficient form of the Bell-polynomial calculus used in the parent article \cite{base}.

\begin{lemma}[Gamma-ratio coefficients]
\label{lem:ratio}
For fixed $a,b$, with $x$ positive and large,
\begin{equation}
 \frac{\Gamma(x+a)}{\Gamma(x+b)}
 \sim x^{a-b}\sum_{n\ge0}R_n(a,b)x^{-n},
 \qquad R_n(a,b)=\cE_n(r(a,b)),
 \label{eq:ratio}
\end{equation}
where
\begin{equation}
 r_j(a,b)=\frac{(-1)^{j+1}}{j(j+1)}
          \bigl(B_{j+1}(a)-B_{j+1}(b)\bigr).
 \label{eq:ratio-log}
\end{equation}
Here $B_m(z)$ denotes the Bernoulli polynomial.
\end{lemma}
\begin{proof}
The shifted Stirling expansion gives
\[
 \log\Gamma(x+a)-\log\Gamma(x+b)
 =(a-b)\log x+\sum_{j=1}^{N}r_j(a,b)x^{-j}
       +\bigO(x^{-N-1}).
\]
The expansion is uniform for $a,b$ in a fixed compact set; see the standard gamma asymptotics and Bernoulli-polynomial conventions in \cite{dlmfgamma,dlmfbernoulli}. Exponentiating a finite truncation and using \eqref{eq:E} proves the formula with a remainder of the indicated order. Since $N$ is arbitrary, it proves the all-orders statement. This use of a divergent formal series does not assert convergence.
\end{proof}

For reference, the Bernoulli convention is fixed by
\[
 \frac{t\ee^{zt}}{\ee^t-1}=\sum_{m\ge0}B_m(z)\frac{t^m}{m!},
 \qquad B_1(z)=z-\tfrac12.
\]
This convention is especially important for the shifted half-integer formulas later in the article.

\subsection{A scaled Lagrange formula}

Suppose an inversion problem has been reduced to
\begin{equation}
 H(x)+R(x)=L,\qquad H(X)=L.
 \label{eq:core-equation}
\end{equation}
The variable $X$ is a \emph{core inverse}. It may already contain a Lambert $W$ function and so absorb infinitely many logarithmic corrections.

\begin{proposition}[Scaled local reversion]
\label{prop:scaled}
Put $x=X(1+z)$ and define the removable quotient
\[
 T_X(u)=\frac{H(X(1+u))-H(X)}{u},\qquad T_X(0)=XH'(X).
\]
Assume $T_X(0)$ is invertible in the coefficient field. In a formal grading in which $R(X(1+u))/T_X(u)$ has positive order, the solution with $z$ of positive order is
\begin{equation}
 z=\sum_{m\ge1}\frac{(-1)^m}{m}
 [u^{m-1}]\left(\frac{R(X(1+u))}{T_X(u)}\right)^m.
 \label{eq:scaledLB}
\end{equation}
For an additive displacement $x=X+z$, the same formula holds with
$T_X(u)=(H(X+u)-H(X))/u$ and $R(X+u)$.
\end{proposition}
\begin{proof}
Equation~\eqref{eq:core-equation} becomes
$z=-R(X(1+z))/T_X(z)$. Insert a bookkeeping parameter $\eta$ multiplying the right-hand side. Ordinary Lagrange--B\"urmann inversion gives the coefficient of $\eta^m$ as the $m$th summand in \eqref{eq:scaledLB}. Positive formal order makes each coefficient in the intended grading a finite sum, so $\eta$ may be set to one formally. Analytically, convergence requires an additional local implicit-function argument; it is not a consequence of this formal calculation alone.
\end{proof}

The point of using a scaled displacement is that the leading inverse correction can be much larger than one. For involutions it is of order $\sqrt X/\log X$; relative to $X$, however, it is small. The wrong choice of an ``infinitesimal'' displacement can obscure an otherwise straightforward reversion.

\subsection{Inverse transport of an exact small sector}
\label{subsec:transport}

A particularly effective approach is to invert the dominant function exactly before expanding its small corrections.

\begin{theorem}[Phase-coordinate sector transport]
\label{thm:transport}
Let $F_0>0$ be $C^\infty$ on an eventual interval with $(\log F_0)^{\prime}>0$, let $\phi_0=\log F_0$, and let $g=\phi_0^{-1}$. Let $\varepsilon$ be $C^\infty$, with $1+\varepsilon>0$, and use a positive sector grading preserved by the differentiations below. Write
\[
 F(x)=F_0(x)(1+\varepsilon(x)),\qquad
 h(L)=\log\bigl(1+\varepsilon(g(L))\bigr).
\]
Then the formal inverse, graded by the small sectors in $\varepsilon$, is
\begin{equation}
 F^{-1}(\ee^L)=g(L)+
 \sum_{m\ge1}\frac{(-1)^m}{m!}
 \left(\frac{\dd}{\dd L}\right)^{m-1}
       \bigl(g'(L)h(L)^m\bigr).
 \label{eq:transport}
\end{equation}
At each fixed locally finite action cutoff, only finitely many sector products occur. The coefficients may contain exact functions of $g(L)$.
\end{theorem}
\begin{proof}
Set $z=\phi_0(x)$. The equation $F(x)=\ee^L$ becomes $z+h(z)=L$. Insert $\eta$ in $z=L-\eta h(z)$ and apply Lagrange--B\"urmann to $g(z)$ rather than to $z$ itself:
\[
 g(z)=g(L)+\sum_{m\ge1}\frac{(-\eta)^m}{m!}
 \left(\frac{\dd}{\dd L}\right)^{m-1}\bigl(g'(L)h(L)^m\bigr).
\]
Setting $\eta=1$ in the sector grading proves the formula. It also shows why derivatives of oscillatory coefficients must not be omitted.
\end{proof}

For example, through second order in an auxiliary sector strength $\eta$, replacing $\varepsilon$ by $\eta\varepsilon$ gives
\begin{align}
 F^{-1}(\ee^L)
 ={}&g-\eta g'\varepsilon(g)\nonumber\\
 &+\frac{\eta^2}{2}\left\{g'\varepsilon(g)^2
       +\frac{\dd}{\dd L}\bigl(g'\varepsilon(g)^2\bigr)\right\}
       +\bigO(\eta^3).
 \label{eq:transport-second}
\end{align}
This displays two sources of nonlinear inverse layers: the logarithm of the forward correction and the change of the evaluation point. Even a missing forward layer can be generated by these operations.

There is an important analytic qualification. An exact $g$ and an exact sector decomposition resolve an exponentially small effect. Replacing $g$ by a fixed finite inverse-power approximation generally does not: the replacement error is usually much larger than that effect. One may either keep $g$ exact, use a rigorously controlled numerical evaluation, or specify a summation and optimal-truncation procedure. We use the first option in the two-saddle and moment examples.

\subsection{Residuals certify inverse errors}

\begin{lemma}[A local inverse certificate]
\label{lem:residual}
Let $f$ be $C^1$ on an interval containing
$[\widehat x-r/m,\widehat x+r/m]$, with $f'(x)\ge m>0$ there. If
$|f(\widehat x)-L|\le r$, then the equation $f(x)=L$ has exactly one solution in that interval and
\begin{equation}
 |x-\widehat x|\le\frac{r}{m}.
 \label{eq:residual}
\end{equation}
\end{lemma}
\begin{proof}
Integrating the derivative lower bound shows that $f(\widehat x-r/m)\le L$ and $f(\widehat x+r/m)\ge L$. Continuity gives existence, strict monotonicity gives uniqueness, and the same derivative bound gives the displayed error.
\end{proof}

For rapidly growing positive functions it is usually best to apply this lemma to $f=\log F$. A relative forward error of size $r$ becomes an inverse error of order $r/(\log F)'$. The relevant slopes in this article are approximately
\[
 \begin{array}{c|c}
 \text{forward growth}&\text{logarithmic slope}\\\hline
 C\ee^{ax}x^\beta&a\\
 (x/\ee)^{x/2}\ee^{\sqrt x}&\tfrac12\log x\\
 \Gamma(x+1)c^x&\log x+\log c\\
 2^{x(x-1)/2}&(x-\tfrac12)\log2
 \end{array}
\]
For harmonic functions, by contrast, the derivative of the function itself is about $1/x$, so forward errors are amplified by a factor of order $x$.

\section{Gamma quotients and the negative Lambert branch}
\label{sec:gamma}

\subsection{A master expansion}

Consider the positive function
\begin{equation}
 F(x)=K\ee^{\mu x}x^\nu
       \prod_{i=1}^r\Gamma(a_ix+b_i)^{\epsilon_i},
 \qquad K>0,\quad a_i>0,
 \label{eq:gamma-general}
\end{equation}
where all parameters are fixed and real. Real powers are defined using the positive gamma values at large positive arguments.

\begin{theorem}[Gamma-quotient coefficient engine]
\label{thm:gamma}
The logarithm of \eqref{eq:gamma-general} has the all-orders expansion
\begin{equation}
 \log F(x)\sim\kappa x\log x+\lambda x+\beta\log x+\log C
                      +\sum_{j\ge1}q_jx^{-j},
 \label{eq:gamma-log}
\end{equation}
where
\begin{align}
 \kappa&=\sum_i\epsilon_i a_i,
 &\lambda&=\mu+\sum_i\epsilon_i a_i(\log a_i-1),\nonumber\\
 \beta&=\nu+\sum_i\epsilon_i(b_i-\tfrac12),
 &\log C&=\log K+\frac{\log(2\pi)}2\sum_i\epsilon_i
           +\sum_i\epsilon_i(b_i-\tfrac12)\log a_i,
 \label{eq:gamma-constants}\\
 q_j&=\frac{(-1)^{j+1}}{j(j+1)}
       \sum_i\frac{\epsilon_iB_{j+1}(b_i)}{a_i^j}.
 \label{eq:gamma-q}
\end{align}
The normalized amplitude coefficients are $\cE_n(q)$.
\end{theorem}
\begin{proof}
Insert the shifted Stirling expansion
\[
 \log\Gamma(ax+b)
 \sim (ax+b-\tfrac12)\log(ax)-ax+\tfrac12\log(2\pi)
 +\sum_{j\ge1}\frac{(-1)^{j+1}B_{j+1}(b)}{j(j+1)(ax)^j}
\]
into the logarithm of \eqref{eq:gamma-general}. Collect $x\log x$, $x$, $\log x$, constants, and inverse powers. There are finitely many gamma factors, so the fixed-order remainders add without changing their order. Formula~\eqref{eq:E} then exponentiates the result. Differentiated versions follow from the corresponding differentiated gamma asymptotics \cite{dlmfgamma}.
\end{proof}

The case $\kappa=0$ is called \emph{balanced}. It includes many combinatorial sequences with purely exponential leading growth multiplied by an algebraic factor. The gamma origin also specifies a real interpolation; it does not specify a unique interpolation among all smooth functions.

\subsection{All-orders inversion of the balanced case}

Assume $\kappa=0$ and $\lambda=a>0$. With $L=\log(y/C)$, define the large solution $X$ of
\begin{equation}
 aX+\beta\log X=L.
 \label{eq:balanced-core}
\end{equation}
Its exact expression is
\begin{equation}
 X=\begin{cases}
 L/a,&\beta=0,\\[1mm]
 \displaystyle\frac{\beta}{a}
 W_{\star}\!\left(\frac a\beta\ee^{L/\beta}\right),&\beta\ne0,
 \end{cases}
 \label{eq:balanced-W}
\end{equation}
where $W_\star=W_0$ for $\beta>0$ and $W_\star=W_{-1}$ for $\beta<0$. The second branch in the negative-$\beta$ case is essential: $W_0$ would give the small rather than the large solution. The branch notation agrees with \texttt{ProductLog[k,z]} \cite{wlproductlog}.

\begin{theorem}[Explicit centered inverse coefficients]
\label{thm:balanced-inverse}
For the specified gamma-quotient interpolation,
\begin{equation}
 F^{-1}(y)\sim X+\sum_{n\ge1}d_nX^{-n}.
 \label{eq:balanced-inverse}
\end{equation}
Let $Q(t)=\sum_{j\ge1}q_jt^j$ and
\[
 \Psi(t,u)=-\frac1a\left\{\beta\log(1+tu)
                     +Q\!\left(\frac{t}{1+tu}\right)\right\}.
\]
Then the nonrecursive finite formula
\begin{equation}
 d_n=\sum_{m=1}^n\frac1m[t^nu^{m-1}]\Psi(t,u)^m
 \label{eq:balanced-d}
\end{equation}
computes every coefficient. In particular,
\begin{align}
 d_1&=-\frac{q_1}{a},\nonumber\\
 d_2&=\frac{\beta q_1}{a^2}-\frac{q_2}{a},\nonumber\\
 d_3&=-\frac{\beta^2q_1}{a^3}
       +\frac{\beta q_2-q_1^2}{a^2}-\frac{q_3}{a}.
 \label{eq:balanced-first}
\end{align}
Truncation through $d_NX^{-N}$ has error $\bigO(X^{-N-1})$.
\end{theorem}
\begin{proof}
Write $x=X+U$, $t=X^{-1}$. Subtracting \eqref{eq:balanced-core} from $\log F(x)=\log y$ gives formally $U=\Psi(t,U)$. Since $\Psi$ has positive $t$-order, Lagrange--B\"urmann gives \eqref{eq:balanced-d}; terms $m>n$ cannot contribute. Direct coefficient comparison yields \eqref{eq:balanced-first}.

For the analytic assertion, substitute a finite truncation into a sufficiently long finite expansion of $\log F$. Coefficient cancellation leaves a residual $\bigO(X^{-N-1})$. Its derivative is $a+\bigO(X^{-1})\ge a/2$ on a fixed neighborhood for large $X$. Lemma~\ref{lem:residual} supplies an inverse with the asserted error and proves its local uniqueness. The derivative asymptotics also prove eventual monotonicity globally.
\end{proof}

Expanding $W_{-1}$ itself converts this compact expression into the familiar nested logarithms. For instance, solving \eqref{eq:balanced-core} by substitution gives
\[
 X=\frac1a\left\{L-\beta\log(L/a)
          +\frac{\beta^2\log(L/a)}{L}
          +\bigO\!\left(\frac{(\log L)^2}{L^2}\right)\right\}.
\]
Keeping $X$ exact avoids repeatedly expanding these logarithmic terms.

\subsection{Catalan numbers}
\label{subsec:catalan}

The Catalan numbers count, among many equivalent objects, rooted full binary plane trees by internal vertices. Their standard formula and gamma continuation are \cite{oeiscatalan,wlcatalan}
\begin{equation}
 C_n=\frac1{n+1}\binom{2n}{n},\qquad
 C(x)=\frac{\Gamma(2x+1)}{\Gamma(x+1)\Gamma(x+2)}.
 \label{eq:catalan}
\end{equation}
They illustrate the distinction between Catalan numbers as coefficients of a quadratic inverse and the inverse of the \emph{counting function} $C(x)$; it is the latter that is studied here.

Theorem~\ref{thm:gamma} gives $a=\log4$, $\beta=-3/2$, $C=\pi^{-1/2}$, and
\begin{align}
 \log C(x)\sim{}&x\log4-\frac32\log x-\frac12\log\pi
 -\frac9{8x}+\frac1{2x^2}-\frac{21}{64x^3}
 +\frac1{4x^4}-\frac{129}{640x^5}+\cdots.
 \label{eq:catalan-log}
\end{align}
Exponentiating,
\begin{equation}
 C(x)\sim\frac{4^x}{\sqrt\pi x^{3/2}}
 \left(1-\frac9{8x}+\frac{145}{128x^2}
        -\frac{1155}{1024x^3}+\frac{36939}{32768x^4}+\cdots\right).
 \label{eq:catalan-amplitude}
\end{equation}
All coefficients are supplied by \eqref{eq:gamma-q} and \eqref{eq:E}, not just the displayed ones.

Define
\begin{equation}
 X=-\frac3{2a}W_{-1}\!\left[
       -\frac{2a}{3}\left(\frac1{\sqrt\pi y}\right)^{2/3}\right],
 \qquad a=\log4.
 \label{eq:catalan-core}
\end{equation}
Then
\begin{align}
 C^{-1}(y)={}&X+\frac9{8aX}
 +\frac1{X^2}\left(\frac{27}{16a^2}-\frac1{2a}\right)\nonumber\\
 &+\frac1{X^3}\left(\frac{81}{32a^3}-\frac{129}{64a^2}
                    +\frac{21}{64a}\right)+\bigO(X^{-4}).
 \label{eq:catalan-inverse}
\end{align}
There is no unknown constant or auxiliary recurrent sequence in this description. Gamma asymptotics and the finite coefficient sum determine every further term.

\subsection{Fuss--Catalan numbers}

For a fixed integer $p\ge2$, full $p$-ary plane trees give the sequence
\begin{equation}
 C_{p,n}=\frac1{(p-1)n+1}\binom{pn}{n},\qquad
 C_p(x)=\frac{\Gamma(px+1)}{\Gamma(x+1)\Gamma((p-1)x+2)}.
 \label{eq:fuss}
\end{equation}
The ternary case is OEIS A001764 \cite{oeisfuss}. The elementary generating-function equation $T=1+zT^p$ and Lagrange inversion also give \eqref{eq:fuss}: for $U=T-1$, $U=z(1+U)^p$, so
$[z^n]U=\binom{pn}{n-1}/n$.

Here
\begin{equation}
 \rho_p=\frac{p^p}{(p-1)^{p-1}},\quad
 a_p=\log\rho_p,\quad \beta=-\frac32,\quad
 C_p^{(0)}=\sqrt{\frac{p}{2\pi(p-1)^3}},
 \label{eq:fuss-constants}
\end{equation}
and
\begin{equation}
 q_j^{(p)}=\frac{(-1)^{j+1}}{j(j+1)}
 \left\{B_{j+1}(1)(p^{-j}-1)
             -\frac{B_{j+1}(2)}{(p-1)^j}\right\}.
 \label{eq:fuss-q}
\end{equation}
For example,
\[
 q_1^{(p)}=\frac1{12}\left(\frac1p-1-\frac{13}{p-1}\right).
\]
Thus $C_p(x)\sim C_p^{(0)}\rho_p^x x^{-3/2}\sum_m\cE_m(q^{(p)})x^{-m}$, and its inverse is obtained from \eqref{eq:balanced-W} and \eqref{eq:balanced-d}. For $p=3$, $\rho_3=27/4$, $C_3^{(0)}=\sqrt3/(4\sqrt\pi)$, and $q_1=-43/72$. This entire fixed-$p$ family shares the same logarithmic exponent, while its growth rate and inverse correction polynomials vary with $p$.

\subsection{Central multinomial coefficients and rectangular tableaux}

For a fixed integer $d\ge2$, consider
\begin{equation}
 M_d(x)=\frac{\Gamma(dx+1)}{\Gamma(x+1)^d}.
 \label{eq:multinomial}
\end{equation}
At $x=n$, this counts assignments of $dn$ labeled objects to $d$ labeled boxes containing $n$ objects each, by the elementary multinomial formula. Its constants are
\begin{equation}
 a=d\log d,\quad \beta=\frac{1-d}{2},\quad
 C=\sqrt d\,(2\pi)^{(1-d)/2}.
 \label{eq:multi-constants}
\end{equation}
Only odd inverse powers occur in its logarithmic correction:
\begin{equation}
 q_{2m-1}=\frac{B_{2m}}{2m(2m-1)}
                 \bigl(d^{1-2m}-d\bigr),\qquad q_{2m}=0.
 \label{eq:multi-q}
\end{equation}
Even inverse powers reappear after exponentiation and inversion. The absence of an order in the logarithm is therefore not the absence of that order in the final inverse.

The number of standard Young tableaux of a rectangle with $d$ rows and $n$ columns is
\begin{equation}
 T_d(n)=(dn)!\frac{\prod_{j=0}^{d-1}j!}{\prod_{j=0}^{d-1}(n+j)!}.
 \label{eq:tableaux}
\end{equation}
One route to this formula is the hook-length formula: the hook lengths in row $i$ are $(n+d-i),(n+d-i-1),\ldots,(d-i+1)$, whose product is $(n+d-i)!/(d-i)!$. Multiplying over rows gives \eqref{eq:tableaux}. The three-row case is recorded as A005789 \cite{oeistableaux}. We take the gamma version of \eqref{eq:tableaux} as the real extension.

Now
\begin{align}
 a&=d\log d,&
 \beta&=-\frac{d^2-1}{2},&
 C&=\sqrt d\,(2\pi)^{(1-d)/2}\prod_{j=0}^{d-1}j!,
 \label{eq:tableaux-constants}\\
 q_r&=\frac{(-1)^{r+1}}{r(r+1)}
 \left\{\frac{B_{r+1}(1)}{d^r}
             -\sum_{j=0}^{d-1}B_{r+1}(j+1)\right\}.
 \label{eq:tableaux-q}
\end{align}
The inverse again uses $W_{-1}$. The case $d=2$ reduces exactly to Catalan numbers, providing a consistency check between two combinatorial descriptions. No assertion here is uniform as $d\to\infty$: varying the rectangle height changes the asymptotic problem.

\subsection{What the gamma engine does not automatically supply}

When $\kappa>0$, the convenient unbalanced core is
\begin{equation}
 X=\frac{L/\kappa}{W_0\!\left((L/\kappa)\ee^{\lambda/\kappa}\right)},
 \qquad L=\log(y/C),
 \label{eq:unbalanced-core}
\end{equation}
which solves $\kappa X\log X+\lambda X=L$. The remaining terms are handled by Proposition~\ref{prop:scaled}. This observation connects the balanced families to the factorial-growth examples below.

However, the Bernoulli expansions of gamma functions are generally divergent. An all-orders expansion is not automatically a complete resurgent transseries with known Stokes multipliers. In this article the gamma-quotient results are rigorous all-orders algebraic sectors and their inverse sectors. We do not attach arbitrary exponentially small terms to them or infer a summation convention from the integer samples.

\section{Stirling numbers: finite exponential grids and nested logarithms}
\label{sec:stirling}

The two kinds of Stirling numbers exhibit remarkably different inverse scales when the second index is fixed. Set partitions produce a finite sum of ordinary exponentials; permutations with a fixed number of cycles produce factorial growth multiplied by a polynomial in $\log n$. The standard combinatorial meanings and sign conventions are described in \cite{dlmfstirling,wlstirling2,wlstirling1}.

\subsection{An exactly solvable multiexponential class}

\begin{theorem}[Convergent multivariate inverse]
\label{thm:finite-exp}
Let
\begin{equation}
 F(x)=A\ee^{ax}\left(1+\sum_{j=1}^r c_j\ee^{-\omega_jx}\right),
 \quad A>0,\quad a>0,\quad \omega_j>0,
 \label{eq:finite-exp}
\end{equation}
with fixed real coefficients. Put
\[
 X=\frac{\log(y/A)}a,\qquad
 \beta_j=\frac{\omega_j}{a},\qquad
 t_j=(y/A)^{-\beta_j}.
\]
For a nonzero multi-index $\bm m$, let $M=\sum_jm_j$ and
$B=\sum_j\beta_jm_j$. Then, for sufficiently large $y$,
\begin{equation}
 F^{-1}(y)=X-\frac1a
 \sum_{\bm m\ne0}
 \frac{(M-1)!}{\prod_jm_j!}
 \binom{B-1}{M-1}\prod_j(c_jt_j)^{m_j}.
 \label{eq:finite-exp-inverse}
\end{equation}
This is an actually convergent multivariate Taylor series near $\bm t=0$, not merely a Poincar\'e expansion. The generalized binomial coefficient is the polynomial
$\binom{z}{m}=z(z-1)\cdots(z-m+1)/m!$.
\end{theorem}
\begin{proof}
Let $\delta=a(x-X)$ and $w=\ee^{-\delta}$. The equation $F(x)=y$ is equivalent to
\begin{equation}
 w=1+\sum_{j=1}^r c_jt_jw^{\beta_j}.
 \label{eq:w-finite}
\end{equation}
The right-hand powers are analytic in $w$ near $1$. At $(w,\bm t)=(1,0)$ the derivative of the left side minus the right side with respect to $w$ is $1$. The analytic implicit-function theorem therefore gives a unique analytic $w(\bm t)$ near that point.

Write $u=w-1$, insert an auxiliary parameter $\eta$ multiplying the sum, and apply Lagrange inversion to $\log(1+u)$:
\[
 \log(1+u)=\sum_{M\ge1}\frac{\eta^M}{M}
 [z^{M-1}]\frac1{1+z}
       \left(\sum_jc_jt_j(1+z)^{\beta_j}\right)^M.
\]
The multinomial coefficient for a given $\bm m$ is $M!/\prod m_j!$, and the remaining extraction is $\binom{B-1}{M-1}$. Since $x=X-a^{-1}\log w$, this proves \eqref{eq:finite-exp-inverse}. The analytic implicit-function construction proves convergence. Finally $F'(x)/(A\ee^{ax})\to a>0$, so the branch is the eventual real inverse.
\end{proof}

When two multi-indices have the same total action $B$, their contributions must be added. A total-degree truncation is useful computationally, but ordering the actual asymptotic terms requires ordering the weights $B$, not just the integers $M$. The action semigroup is locally finite because all $\beta_j$ are positive and there are finitely many generators.

\subsection{Stirling numbers of the second kind and surjections}

For a fixed integer $k\ge2$, inclusion--exclusion gives
\begin{equation}
 S_k(x)=\frac1{k!}\sum_{j=1}^k(-1)^{k-j}\binom{k}{j}j^x,
 \qquad x>0.
 \label{eq:stirling2-cont}
\end{equation}
At positive integers $n$, this equals $\StirII nk$. Indeed, among the $k^n$ maps from an $n$-element labeled set to $k$ labeled boxes, inclusion--exclusion removes maps omitting one or more boxes. Dividing the number of surjections by $k!$ forgets the labels of the nonempty boxes. Equation~\eqref{eq:stirling2-cont} is also our exact continuous interpolation.

Theorem~\ref{thm:finite-exp} applies with
\begin{equation}
 A=\frac1{k!},\qquad a=\log k,\qquad
 c_j=(-1)^{k-j}\binom{k}{j},\qquad
 \omega_j=\log\frac{k}{j}\quad(1\le j<k).
 \label{eq:stirling2-params}
\end{equation}
Thus every inverse coefficient is given by the finite generalized-binomial expression in \eqref{eq:finite-exp-inverse}. Multiplying the forward function by $k!$ gives the corresponding inverse for the surjection count, simply by replacing the target $y$ by $y/k!$.

For $k=2$ the formula is elementary:
\[
 S_2(x)=2^{x-1}-1,\qquad
 S_2^{-1}(y)=1+\frac{\log(y+1)}{\log2}.
\]
Its convergent expansion in $1/y$ provides a check on the general theorem.

For $k=3$, define
\[
 \beta=\frac{\log(3/2)}{\log3},\qquad
 X=\frac{\log(6y)}{\log3},\qquad
 t=(6y)^{-\beta},\quad s=(6y)^{-1}.
\]
Since $S_3(x)=(3^x-3\cdot2^x+3)/6$, the inverse starts as
\begin{align}
 (\log3)(S_3^{-1}(y)-X)
 ={}&3t-3s+\frac92(1-2\beta)t^2+9\beta ts-\frac92s^2
 \nonumber\\
 &+\bigO((|t|+|s|)^3).
 \label{eq:stirling3-inverse}
\end{align}
This display is arranged by multivariate degree. In dominance order, $t^2$ precedes $s$, since $2\beta<1$. The pure cubic term, useful for a dominance-ordered expansion, is
\[
 \frac92(3\beta-1)(3\beta-2)t^3.
\]
The complete formula automatically places mixed terms and accounts for any collisions. Unlike a factorial asymptotic series, this inverse expansion genuinely converges for sufficiently large target values.

\subsection{Unsigned Stirling numbers of the first kind}

Write $c(n,k)=\StirI nk$ for the number of permutations of $n$ labels with $k$ cycles. The polynomial identity
\[
 z(z+1)\cdots(z+n-1)=\sum_{k=0}^n c(n,k)z^k
\]
follows by adding the new label either as a singleton cycle or in one of the existing cycle positions. It gives a particularly useful interpolation for each fixed $k\ge1$:
\begin{equation}
 C_k(x)=[z^k]\frac{\Gamma(x+z)}{\Gamma(z)},\qquad x>0.
 \label{eq:stirling1-cont}
\end{equation}
The coefficient is taken in the analytic germ at $z=0$; $1/\Gamma(z)$ has a zero there. This is an explicit function, not a recurrence.

The Wolfram Language function \texttt{StirlingS1[n,k]} uses the \emph{signed} convention, so at integers
\[
 c(n,k)=(-1)^{n-k}\,\texttt{StirlingS1[n,k]}.
\]
The gamma expression \eqref{eq:stirling1-cont} fixes the positive unsigned continuation; analytically continuing $(-1)^{x-k}$ would produce a different object.

Using the Taylor expansions of $\log\Gamma$ at $x$ and at $1$ gives the exact finite-coefficient representation
\begin{align}
 C_k(x)=\Gamma(x)[z^{k-1}]\exp\Bigg(& (\psi(x)+\gamma)z\nonumber\\
 &+\sum_{j\ge2}\left\{\frac{\psi^{(j-1)}(x)}{j!}
                  -\frac{(-1)^j\zeta(j)}j\right\}z^j\Bigg).
 \label{eq:stirling1-polygamma}
\end{align}
Only $j\le k-1$ can affect this coefficient. Here $\psi$ is the logarithmic derivative of gamma and $\gamma$ is Euler's constant.

\begin{theorem}[All-orders fixed-column expansion]
\label{thm:stirling1}
Put $\ell=\log x$ and
\begin{align}
 r_j(z)&=\frac{(-1)^{j+1}}{j(j+1)}
              \bigl(B_{j+1}(z)-B_{j+1}(0)\bigr),\nonumber\\
 P_{k,m}(\ell)&=[z^{k-1}]\frac{\ee^{\ell z}}{\Gamma(1+z)}
                      \cE_m(r(z)).
 \label{eq:stirling1-P}
\end{align}
For fixed $k$,
\begin{equation}
 C_k(x)\sim\Gamma(x)\sum_{m\ge0}P_{k,m}(\log x)x^{-m}.
 \label{eq:stirling1-forward}
\end{equation}
The leading polynomial $P_{k,0}$ has degree $k-1$, and every $P_{k,m}$ with $m\ge1$ has degree at most $k-2$. In particular,
\begin{align}
 P_{1,0}&=1,&P_{2,0}&=\ell+\gamma,\nonumber\\
 P_{3,0}&=\frac{(\ell+\gamma)^2-\zeta(2)}2,&
 P_{4,0}&=\frac{(\ell+\gamma)^3-3\zeta(2)(\ell+\gamma)+2\zeta(3)}6.
 \label{eq:stirling1-leading}
\end{align}
\end{theorem}
\begin{proof}
Rewrite \eqref{eq:stirling1-cont} as
\[
 C_k(x)=\Gamma(x)[z^{k-1}]
        \frac1{\Gamma(1+z)}\frac{\Gamma(x+z)}{\Gamma(x)}.
\]
Lemma~\ref{lem:ratio} expands the gamma ratio uniformly on a fixed small circle about $z=0$. Cauchy's coefficient formula therefore permits extraction term by term, with a fixed-order remainder bounded by an inverse power times a fixed power of $\log x$. The expression \eqref{eq:stirling1-P} follows.

For $m\ge1$, the polynomial $\cE_m(r(z))$ vanishes at $z=0$, because each $r_j(0)=0$. Its product with $\ee^{\ell z}$ can consequently contribute at most $\ell^{k-2}$ to $[z^{k-1}]$. For $m=0$ the leading term is $\ell^{k-1}/(k-1)!$. Finally,
\[
 \frac1{\Gamma(1+z)}
 =\exp\left(\gamma z+\sum_{j\ge2}\frac{(-1)^{j+1}\zeta(j)}jz^j\right)
\]
gives the displayed leading polynomials by finite extraction.
\end{proof}

\subsection{The nested-logarithmic inverse}

Define
\begin{equation}
 L=\log\left(\frac{y(k-1)!}{\sqrt{2\pi}}\right),\qquad
 X=\frac{L}{W_0(L/\ee)},\qquad \ell=\log X,
 \label{eq:stirling1-core}
\end{equation}
and the polynomial in $1/\ell$
\begin{equation}
 Q_k(1/\ell)=\frac{(k-1)!P_{k,0}(\ell)}{\ell^{k-1}}.
 \label{eq:stirling1-Q}
\end{equation}
It tends to one. Taking logarithms in Theorem~\ref{thm:stirling1}, and also using Stirling's formula for $\Gamma(x)$, produces
\begin{align}
 \log\left(\frac{C_k(x)(k-1)!}{\sqrt{2\pi}}\right)
 \sim{}&x\log x-x-\tfrac12\log x
       +\log\bigl((k-1)!P_{k,0}(\log x)\bigr)
 \nonumber\\
 &+\sum_{m\ge1}U_{k,m}(\log x)x^{-m}.
 \label{eq:stirling1-log}
\end{align}
Each $U_{k,m}$ is explicit: add the gamma logarithmic coefficient at order $m$ to
\[
 [t^m]\log\left(1+\sum_{j\ge1}
       \frac{P_{k,j}(\ell)}{P_{k,0}(\ell)}t^j\right).
\]
Thus the coefficients are rational functions of $\ell$, which themselves possess expansions in $1/\ell$.

\begin{corollary}[First inverse layers]
\label{cor:stirling1-inverse}
For the interpolation \eqref{eq:stirling1-cont},
\begin{equation}
 C_k^{-1}(y)=X+\frac12
       -\frac{(k-1)\log\ell}{\ell}
       -\frac{\log Q_k(1/\ell)}{\ell}
       +\bigO(X^{-1}).
 \label{eq:stirling1-inverse}
\end{equation}
In particular,
\begin{equation}
 C_k^{-1}(y)=X+\frac12-
       \frac{(k-1)\log\log X}{\log X}
       -\frac{(k-1)\gamma}{(\log X)^2}
       +\bigO((\log X)^{-3}).
 \label{eq:stirling1-inverse-expanded}
\end{equation}
All further $X^{-1}$ and logarithmic layers are obtained by applying Proposition~\ref{prop:scaled} to \eqref{eq:stirling1-log}.
\end{corollary}
\begin{proof}
The core $H(x)=x\log x-x$ has $H'(X)=\ell$. Its residual at $X$ is
$-\ell/2+(k-1)\log\ell+\log Q_k(1/\ell)+\bigO(X^{-1})$.
Subtracting this residual divided by $\ell$ gives \eqref{eq:stirling1-inverse}. The proposed displacement is bounded. Taylor's theorem then leaves a residual $\bigO(X^{-1})$, whose inverse error is no larger than $\bigO(X^{-1})$ because the slope grows like $\ell$. Also
$Q_k(1/\ell)=1+(k-1)\gamma/\ell+\bigO(\ell^{-2})$, proving the second display. Differentiated asymptotics show positivity and eventual monotonicity, so the continuous inverse exists.
\end{proof}

A useful practical distinction emerges. For fixed $k$, second-kind inverses are logarithmic cores plus convergent generalized powers of $1/y$. First-kind inverses have a Lambert core, an order-one shift, and a hierarchy involving $\log\log X/\log X$ before any $X^{-1}$ correction is reached. The two ``Stirling'' inverse problems belong to different scales.

\section{Moment sequences and their endpoint sectors}
\label{sec:moments}

\subsection{An endpoint coefficient lemma}

A recurring integral has, near its largest positive support point $b$, the form
\[
 K\int t^{x+\alpha-1}(b-t)^{\nu-1}A(b-t)\dd t,
 \qquad \nu>0,
\]
where $A$ is analytic at zero. The change $t=b(1-v)$ reduces the endpoint to a beta-type integral. Taylor expansion of $A(bv)$ and Lemma~\ref{lem:ratio} then produce every inverse-power coefficient.

\begin{lemma}[Endpoint extraction]
\label{lem:endpoint}
Suppose the part of an integral away from $v=0$ is exponentially smaller than its $v=0$ endpoint contribution, and its local integrand is
\[
 Kb^x(1-v)^{x+\alpha-1}v^{\nu-1}(1-\kappa v)^\theta.
\]
After normalizing by $Kb^x\Gamma(\nu)x^{-\nu}$, its inverse-power coefficient at order $m$ is
\begin{equation}
 a_m=\sum_{j=0}^m\binom\theta j(-\kappa)^j
          \rising\nu j\,R_{m-j}(\alpha,\alpha+\nu+j).
 \label{eq:endpoint-coeff}
\end{equation}
\end{lemma}
\begin{proof}
On a small fixed interval about zero, expand $(1-\kappa v)^\theta$ through order $N$. The remainder is bounded by a constant times $v^{N+1}$. Its integral is $\bigO(x^{-\nu-N-1})$, as follows either by comparison with $\ee^{-xv}$ or by a beta integral. Extending each polynomial term to $v=1$ changes it only exponentially, since the extension begins a fixed distance from zero. The $j$th beta integral is
\[
 B(x+\alpha,\nu+j)=\Gamma(\nu+j)
 \frac{\Gamma(x+\alpha)}{\Gamma(x+\alpha+\nu+j)}.
\]
Apply Lemma~\ref{lem:ratio}, divide by $\Gamma(\nu)x^{-\nu}$, and collect $j+(m-j)=m$. This proves both the coefficient formula and the fixed-order remainder. It is the endpoint form of Watson's lemma \cite{dlmflaplace}.
\end{proof}

The proof is local. It does not assert convergence of an infinite endpoint Taylor expansion throughout the whole integration interval. In the Motzkin example we have a stronger globally convergent beta representation; the Delannoy and Schr\"oder applications only need this local lemma.

\subsection{Motzkin numbers as two exact moment sectors}
\label{subsec:motzkin}

The Motzkin numbers $1,1,2,4,9,\ldots$ count paths with up, down, and horizontal steps that stay above the horizontal axis; their generating function is recorded in A001006 \cite{oeismotzkin}. They have the exact moment representation
\begin{equation}
 M_n=\frac1{2\pi}\int_{-1}^3t^n\sqrt{(3-t)(t+1)}\dd t.
 \label{eq:motzkin-moment}
\end{equation}
To verify it directly, put $t=1+u$. The centered semicircle moments of odd order vanish, and for even order $2j$ the substitution $u=2\cos\theta$ gives
\[
 \frac1{2\pi}\int_{-2}^2u^{2j}\sqrt{4-u^2}\dd u
 =\frac1{j+1}\binom{2j}{j}=C_j.
\]
Therefore the right side of \eqref{eq:motzkin-moment} equals
$\sum_j\binom n{2j}C_j$. Its generating function is
\[
 \frac1{1-z}\sum_{j\ge0}C_j\left(\frac{z^2}{(1-z)^2}\right)^j
 =\frac{1-z-\sqrt{1-2z-3z^2}}{2z^2},
\]
which proves the identity.

Define positive functions for $x>-1$ by
\begin{align}
 M_+(x)&=\frac1{2\pi}\int_0^3t^x\sqrt{(3-t)(t+1)}\dd t,\nonumber\\
 M_-(x)&=\frac1{2\pi}\int_0^1u^x\sqrt{(3+u)(1-u)}\dd u.
 \label{eq:motzkin-sectors}
\end{align}
Then
\begin{equation}
 M(x)=M_+(x)+\cos(\pi x)M_-(x)
 \label{eq:motzkin-cont}
\end{equation}
is a specified real interpolation of $M_n$. The cosine is not forced by the integer data, but once chosen it fixes the subdominant real continuation exactly.

\begin{proposition}[Convergent beta decompositions]
\label{prop:motzkin-beta}
For $x>-1$,
\begin{align}
 M_+(x)&=\frac{3\sqrt3}{\pi}\,3^x
    \sum_{j\ge0}\binom{1/2}{j}\left(-\frac34\right)^j
          B(x+1,j+\tfrac32),\nonumber\\
 M_-(x)&=\frac1\pi
    \sum_{j\ge0}\binom{1/2}{j}\left(-\frac14\right)^j
          B(x+1,j+\tfrac32).
 \label{eq:motzkin-beta}
\end{align}
Both series converge absolutely.
\end{proposition}
\begin{proof}
For the first integral use $t=3(1-v)$, and for the second use $u=1-v$. The remaining square roots are respectively $(1-3v/4)^{1/2}$ and $(1-v/4)^{1/2}$. Their binomial series converge absolutely and uniformly for $0\le v\le1$. The other factors are integrable for $x>-1$, so termwise integration is justified and gives the beta functions in \eqref{eq:motzkin-beta}.
\end{proof}

Applying Lemma~\ref{lem:endpoint} to each sector gives
\begin{align}
 M_+(x)&\sim\frac{3\sqrt3}{2\sqrt\pi}\,3^xx^{-3/2}
                         \sum_{m\ge0}a_m^+x^{-m},\nonumber\\
 M_-(x)&\sim\frac1{2\sqrt\pi}\,x^{-3/2}
                         \sum_{m\ge0}a_m^-x^{-m},
 \label{eq:motzkin-asymp}
\end{align}
with the finite formulas
\begin{equation}
 a_m^\pm=\sum_{j=0}^m\binom{1/2}{j}(-\kappa_\pm)^j
       \rising{3/2}{j}R_{m-j}(1,j+\tfrac52),
 \quad \kappa_+=\tfrac34,\quad\kappa_-=\tfrac14.
 \label{eq:motzkin-coeff}
\end{equation}
The first coefficients are
\begin{center}
\begin{tabular}{c@{\qquad}rrrrr}
\toprule
 & $m=0$&$m=1$&$m=2$&$m=3$&$m=4$\\\midrule
$a_m^+$&$1$&$-39/16$&$2665/512$&$-87885/8192$&$11422299/524288$\\[2pt]
$a_m^-$&$1$&$-33/16$&$1945/512$&$-57435/8192$&$6977019/524288$\\\bottomrule
\end{tabular}
\end{center}
In particular,
\begin{equation}
 \frac{M_-(x)}{M_+(x)}
 \sim\frac{3^{-x}}{3\sqrt3}\left(1+\frac3{8x}+\cdots\right).
 \label{eq:motzkin-ratio}
\end{equation}
This is a genuine, exactly normalized second sector, not a guessed correction inferred from a recurrence.

\subsection{The Motzkin inverse, including the alternating sector}

For the dominant sector, Theorem~\ref{thm:balanced-inverse} applies with
\[
 a=\log3,\quad\beta=-\frac32,\quad C=\frac{3\sqrt3}{2\sqrt\pi},
 \qquad q_j=\cL_j(a^+).
\]
Although $M_+$ is not itself a finite gamma quotient, the proof of that theorem uses only its differentiated all-orders expansion and a positive limiting logarithmic slope. The moment integral supplies these hypotheses. For instance,
\[
 q_1=-\frac{39}{16},\qquad q_2=\frac{143}{64},\qquad
 M_+^{-1}(y)=X+\frac{39}{16aX}
  +\frac1{X^2}\left(\frac{117}{32a^2}-\frac{143}{64a}\right)+\cdots,
\]
where $aX-\tfrac32\log X=\log(y/C)$.

For the small sector, retain the \emph{exact} dominant inverse $x_+=M_+^{-1}(y)$. Taylor expansion of \eqref{eq:motzkin-cont} and a derivative bound give
\begin{equation}
 M^{-1}(y)=x_+-\frac{\cos(\pi x_+)M_-(x_+)}{M_+'(x_+)}
                +\bigO(3^{-2x_+}).
 \label{eq:motzkin-inverse-exact}
\end{equation}
Equivalently, with a uniform additive error even at zeros of the cosine,
\begin{equation}
 M^{-1}(y)=x_+-\frac{\cos(\pi x_+)}{3\sqrt3\log3}\,3^{-x_+}
 +\bigO\!\left(\frac{3^{-x_+}}{x_+}\right)
 +\bigO(3^{-2x_+}).
 \label{eq:motzkin-inverse-leading}
\end{equation}
All higher exponential layers and their trigonometric coefficients follow from Theorem~\ref{thm:transport} with
$\varepsilon(x)=\cos(\pi x)M_-(x)/M_+(x)$. The exact integrals and the convergent beta representations make the meaning of these layers unambiguous. Differentiated endpoint asymptotics imply $M_+'(x)/M_+(x)\to\log3$, while the other term and its derivative are exponentially smaller; hence $M$ is eventually increasing.

\subsection{Central Delannoy numbers}

Set
\begin{equation}
 a_0=3-2\sqrt2,\qquad b_0=3+2\sqrt2,\qquad
 d=b_0-a_0=4\sqrt2,\qquad \kappa=\frac{b_0}{d}.
 \label{eq:DS-constants}
\end{equation}
The central Delannoy numbers, A001850, have generating function
$(1-6z+z^2)^{-1/2}$ \cite{oeisdelannoy}. A positive real interpolation is
\begin{align}
 D(x)&=\frac1\pi\int_0^\pi(3+2\sqrt2\cos\theta)^x\dd\theta\nonumber\\
 &=\frac1\pi\int_{a_0}^{b_0}
           \frac{t^x}{\sqrt{(b_0-t)(t-a_0)}}\dd t.
 \label{eq:delannoy-moment}
\end{align}
It is also $P_x(3)$ in the usual Legendre-function notation \cite{dlmflegendre}. To check the integer sequence directly, summing the geometric series in the first integral and substituting $u=\tan(\theta/2)$ gives
\[
 \frac1\pi\int_0^\pi\frac{\dd\theta}{1-z(3+2\sqrt2\cos\theta)}
 =\frac1{\sqrt{(1-a_0z)(1-b_0z)}}.
\]

The endpoint $t=b_0$ yields
\begin{equation}
 D(x)\sim C_D b_0^x x^{-1/2}\sum_{m\ge0}d_mx^{-m},
 \qquad C_D=\sqrt{\frac{b_0}{\pi d}},
 \label{eq:delannoy-asymp}
\end{equation}
where
\begin{equation}
 d_m=\sum_{j=0}^m\binom{-1/2}{j}(-\kappa)^j
             \rising{1/2}{j}R_{m-j}(1,j+\tfrac32).
 \label{eq:delannoy-coeff}
\end{equation}
For example,
\[
 d_0=1,\qquad d_1=-\frac14+\frac{3\sqrt2}{32},\qquad
 d_2=\frac{113}{1024}-\frac{9\sqrt2}{128}.
\]
The inverse has the balanced core with $a=\log b_0$, $\beta=-1/2$, $C=C_D$, and coefficients $q_j=\cL_j(d)$. Its first displacement is
\[
 D^{-1}(y)=X+\frac{1/4-3\sqrt2/32}{(\log b_0)X}+\bigO(X^{-2}).
\]
All orders follow from the two finite coefficient formulas \eqref{eq:delannoy-coeff} and \eqref{eq:balanced-d}.

\subsection{Large Schr\"oder numbers}

The large Schr\"oder numbers $1,2,6,22,90,\ldots$, A006318, have generating function
$(1-z-\sqrt{1-6z+z^2})/(2z)$ \cite{oeisschroder}. Define
\begin{equation}
 S(x)=\frac1{2\pi}\int_{a_0}^{b_0}
                t^{x-1}\sqrt{(b_0-t)(t-a_0)}\dd t.
 \label{eq:schroder-moment}
\end{equation}
This integral exists for every real $x$ because $a_0>0$. Its mass at $x=0$ is one: the elementary integral of the square-root density divided by $t$ is
$\pi((a_0+b_0)/2-\sqrt{a_0b_0})=2\pi$. Applying the same elementary Stieltjes-transform calculation as above to
$1/[t(1-zt)]=1/t+z/(1-zt)$ gives the stated generating function. Thus $S(n)$ is the desired integer sequence.

Endpoint extraction now yields
\begin{equation}
 S(x)\sim C_S b_0^xx^{-3/2}\sum_{m\ge0}s_mx^{-m},
 \qquad C_S=\frac{\sqrt{b_0d}}{4\sqrt\pi},
 \label{eq:schroder-asymp}
\end{equation}
with
\begin{equation}
 s_m=\sum_{j=0}^m\binom{1/2}{j}(-\kappa)^j
                  \rising{3/2}{j}R_{m-j}(0,j+\tfrac32).
 \label{eq:schroder-coeff}
\end{equation}
The shift zero in the gamma ratio is significant: the moment integrand is $t^{x-1}$, not $t^x$. The first coefficients are
\[
 s_0=1,\qquad s_1=-\frac34-\frac{9\sqrt2}{32},\qquad
 s_2=\frac{665}{1024}+\frac{45\sqrt2}{128}.
\]
Again the inverse is given to all orders by $a=\log b_0$, $\beta=-3/2$, $C=C_S$, and $q_j=\cL_j(s)$. In particular,
\[
 S^{-1}(y)=X+\frac{3/4+9\sqrt2/32}{(\log b_0)X}+\bigO(X^{-2}).
\]

\subsection{A boundary on the endpoint interpretation}

Both $D(x)$ and $S(x)$ have exact positive-support moment definitions. These definitions fix their real inverses and provide differentiated all-orders expansions. They do \emph{not}, without further analysis, justify adding a nonzero lower-endpoint exponential sector merely because the generating function has a second algebraic singularity or a recurrence has a second characteristic root. A solution of the recurrence can have a zero coefficient in that formal mode, and endpoint decompositions can depend on a contour or a summation convention.

Accordingly, the Delannoy and Schr\"oder results above claim the full algebraic sector and its inverse, not a computed hyperasymptotic decomposition. Motzkin numbers are different: splitting an actual real moment integral at zero produced two exact functions before any expansion was made. That stronger input is what licensed the explicit exponentially small inverse term.

\section{Involutions: two real saddles and a stretched-exponential inverse}
\label{sec:involutions}

\subsection{An exact Gaussian decomposition}

An involution consists of fixed points and disjoint transpositions. Choosing $j$ transpositions gives
\begin{equation}
 I_n=\sum_{j=0}^{\lfloor n/2\rfloor}
        \frac{n!}{2^jj!(n-2j)!},\qquad
 \sum_{n\ge0}I_n\frac{z^n}{n!}=\ee^{z+z^2/2}.
 \label{eq:involution-enumeration}
\end{equation}
These are the involution or telephone numbers, A000085 \cite{oeisinvolution}. If $Z$ is a standard normal random variable, its exponential moment gives $I_n=\mathbb E(1+Z)^n$. Equivalently, define for $x>-1$
\begin{equation}
 A_\sigma(x)=\frac1{\sqrt{2\pi}}\int_0^\infty
          t^x\exp\!\left(-\frac{(t-\sigma)^2}{2}\right)\dd t,
 \qquad \sigma\in\{+1,-1\}.
 \label{eq:involution-A}
\end{equation}
Splitting the Gaussian moment at zero gives
\begin{equation}
 I_n=A_+(n)+(-1)^nA_-(n).
 \label{eq:involution-exact-lattice}
\end{equation}
We select the real interpolation
\begin{equation}
 I(x)=A_+(x)+\cos(\pi x)A_-(x).
 \label{eq:involution-cont}
\end{equation}
Both $A_\sigma$ are positive exact functions. They are therefore suitable sector normalizations: no subdominant constant is being selected by truncating a divergent dominant series.

\subsection{The moving saddle}

For fixed $x$, the logarithm of the integrand in \eqref{eq:involution-A} is
\[
 h_\sigma(t)=x\log t-\frac{(t-\sigma)^2}{2}.
\]
Its unique maximum occurs at
\begin{equation}
 r_\sigma=\frac{\sigma+\sqrt{1+4x}}2,
 \qquad r_\sigma^2-\sigma r_\sigma=x,
 \label{eq:involution-saddle}
\end{equation}
and its curvature is
\begin{equation}
 b_\sigma=-h_\sigma''(r_\sigma)
          =2-\frac{\sigma}{r_\sigma}.
 \label{eq:involution-curvature}
\end{equation}
In particular, the saddle moves like $\sqrt x$, but its Gaussian width remains of order one.

Expand at $r=r_\sigma$ using $u=t-r$. For $k\ge3$,
\[
 \frac{h_\sigma^{(k)}(r)}{k!}
       =\frac{(-1)^{k-1}x}{kr^k}.
\]
Thus the phase has the form
\[
 h_\sigma(r)-\frac{b_\sigma u^2}{2}
       +\sum_{k\ge3}\frac{(-1)^{k-1}x}{kr^k}u^k.
\]
Multiplying these Taylor terms and integrating Gaussian moments is enough to compute every coefficient. The following theorem turns that procedure into a single finite coefficient formula.

\begin{theorem}[All coefficients of both involution saddles]
\label{thm:involution-coeff}
Put $t=x^{-1/2}$ and, in the formulas below, use $t$ as a formal variable. Define
\begin{equation}
 R_\sigma(t)=\frac{\sigma t+\sqrt{4+t^2}}2,
 \qquad b_\sigma(t)=2-\frac{\sigma t}{R_\sigma(t)}.
 \label{eq:involution-R}
\end{equation}
For a finitely supported vector $(m_3,m_4,\ldots)$ of nonnegative integers, set
$K=\sum_{k\ge3}km_k$ and $M=\sum_{k\ge3}m_k$. Define the formal Gaussian correction
\begin{equation}
 \mathcal S_\sigma(t)=
 \sum_{\substack{(m_k)\ge0\\K\ \mathrm{even}}}
 \frac{(-1)^M(K-1)!!}{\prod_{k\ge3}m_k!k^{m_k}}
 \frac{t^{K-2M}}{R_\sigma(t)^K b_\sigma(t)^{K/2}},
 \label{eq:involution-S}
\end{equation}
where the empty vector contributes $1$ and $(-1)!!=1$. Finally put
\begin{align}
 B_\sigma(t)={}&
 \exp\!\left(\frac{\log R_\sigma(t)}{t^2}
       +\frac{\sigma R_\sigma(t)}{2t}-\frac\sigma t-\frac14\right)
 \sqrt{\frac2{b_\sigma(t)}}\,\mathcal S_\sigma(t).
 \label{eq:involution-B}
\end{align}
The apparent negative powers in the exponent cancel. Every coefficient $[t^N]B_\sigma(t)$ is a finite sum: only vectors with $K-2M\le N$ are needed. If $B(t)=B_+(t)$, then $B_-(t)=B(-t)$, and
\begin{equation}
 A_\sigma(x)\sim
 \frac{\ee^{-1/4}}{\sqrt2}\left(\frac{x}{\ee}\right)^{x/2}
          \ee^{\sigma\sqrt x}B(\sigma x^{-1/2}).
 \label{eq:involution-saddle-expansion}
\end{equation}
\end{theorem}

\begin{proof}
The normalized Gaussian moment of $u^K$ is zero for odd $K$ and equals $(K-1)!!\,b_\sigma^{-K/2}$ for even $K$. For a vector $(m_k)$, the product of Taylor coefficients contributes
\[
 \prod_{k\ge3}\frac1{m_k!}
       \left(\frac{(-1)^{k-1}x}{kr^k}\right)^{m_k}.
\]
When $K$ is even, its sign is $(-1)^{K-M}=(-1)^M$. Since $x=t^{-2}$ and $r=R_\sigma/t$, this gives exactly \eqref{eq:involution-S}. Each nonzero $m_k$ contributes positive weight $k-2$, so extraction at a fixed power is finite.

The saddle identity $r^2-\sigma r=x$ simplifies the phase to
\[
 h_\sigma(r)=x\log r-\frac x2+\frac{\sigma r}{2}-\frac12.
\]
Dividing $\ee^{h_\sigma(r)}/\sqrt{b_\sigma}$ by the leading factor in \eqref{eq:involution-saddle-expansion} yields the exponential and square root in \eqref{eq:involution-B}. The expansions $R_\sigma=1+\sigma t/2+t^2/8+\cdots$ and $\log R_\sigma=\sigma\operatorname{arsinh}(t/2)$ show cancellation of the negative powers. Replacing $(\sigma,t)$ by $(+,\sigma t)$ proves the sign symmetry; $K-2M$ is even whenever $K$ is even.

For a rigorous fixed-order remainder, note that
$h_\sigma''(v)=-x/v^2-1\le-1$ for $v>0$. The phase is globally strictly concave, and the contribution outside $|v-r_\sigma|\le x^\epsilon$ is bounded relative to its maximum by a Gaussian tail $\bigO(\ee^{-c x^{2\epsilon}})$. On the inner interval, Taylor expansion to a sufficiently high finite order is uniform; choose $\epsilon>0$ small relative to the desired order. Its remainder is a polynomially bounded factor times a higher power of $x^{-1/2}$, integrable against a slightly weaker Gaussian. Extending the local Gaussian integral to the real line changes it by another negligible tail. This proves the expansion to every fixed order. The same argument after finitely many differentiations supplies the differentiated expansions used in inversion. This is a moving-saddle version of Laplace's method \cite{dlmflaplace}.
\end{proof}

The first terms are
\begin{align}
 B(t)={}&1+\frac7{24}t-\frac{119}{1152}t^2
       -\frac{7933}{414720}t^3
       +\frac{1967381}{39813120}t^4
       -\frac{57200419}{1337720832}t^5+\cdots,
 \label{eq:involution-B-first}\\
 \log B(t)={}&\frac7{24}t-\frac7{48}t^2
      +\frac{37}{1920}t^3+\frac5{128}t^4
      -\frac{5879}{107520}t^5+\cdots.
 \label{eq:involution-logB}
\end{align}
The amplitude coefficients in \eqref{eq:involution-B-first} agree with the independent high-order expansion of Ekhad, Kauers, and Zeilberger \cite{ekz}. They are included here as checks of the general saddle formula, not as a claim of priority for those numerical coefficients.

Combining the exact decomposition and the two expansions gives
\begin{equation}
 I(x)\sectorsim
 \frac{\ee^{-1/4}}{\sqrt2}\left(\frac{x}{\ee}\right)^{x/2}\ee^{\sqrt x}
 \left\{B(x^{-1/2})+\cos(\pi x)\ee^{-2\sqrt x}B(-x^{-1/2})\right\}.
 \label{eq:involution-transseries}
\end{equation}
The second sector is stretched-exponential relative to the first. At integers its sign alternates. A finite truncation of $B$ has an error much larger than that second sector; the exact functions $A_\pm$ retain its meaning despite this fact.

\subsection{Half-power and logarithmic layers of the inverse}

First invert $A_+$, and write $x_+=A_+^{-1}(y)$. Define
\begin{equation}
 L=\log(y\sqrt2)+\frac14,\qquad
 X=\frac{2L}{W_0(2L/\ee)},\qquad \ell=\log X.
 \label{eq:involution-core}
\end{equation}
The core solves $X(\log X-1)/2=L$. Its omission of the $\sqrt x$ term creates an inverse displacement of order $\sqrt X/\log X$, not merely an inverse-power displacement.

\begin{theorem}[Involution dominant inverse]
\label{thm:involution-inverse}
The inverse of the positive saddle has an all-orders expansion
\begin{equation}
 x_+\sim X+\sum_{j\ge0}z_j(\ell)X^{(1-j)/2},
 \label{eq:involution-inverse-general}
\end{equation}
where every $z_j$ is a rational function of $\ell$. The first three are
\begin{align}
 z_0(\ell)&=-\frac2\ell,\nonumber\\
 z_1(\ell)&=\frac2{\ell^2}-\frac2{\ell^3},\nonumber\\
 z_2(\ell)&=-\frac7{12\ell}-\frac1{\ell^3}
                +\frac{14}{3\ell^4}-\frac4{\ell^5}.
 \label{eq:involution-z}
\end{align}
In particular, using these three terms leaves an error $\bigO(X^{-1})$.
\end{theorem}
\begin{proof}
Write $t=X^{-1/2}$, $x=X(1+tz)$, and $\log B(t)=\sum_{j\ge1}q_jt^j$. Subtract the core equation from $\log A_+(x)=\log y$, and divide by $\sqrt X$. The resulting formal equation is
\begin{align}
 0={}&\frac\ell2z+
 \frac{(1+tz)\log(1+tz)-tz}{2t}
 +\sqrt{1+tz}
 +\sum_{j\ge1}q_jt^{j+1}(1+tz)^{-j/2}.
 \label{eq:involution-z-equation}
\end{align}
At $t=0$ this reduces to $\ell z/2+1=0$, so $z_0=-2/\ell$. The derivative with respect to $z$ is $\ell/2$, an invertible element of $\mathbb Q(\ell)$. The formal implicit-function theorem therefore gives a unique $z(t)=\sum_{j\ge0}z_j(\ell)t^j$. Comparing the first three coefficients gives \eqref{eq:involution-z}.

For a finite, nonrecursive formula, let the sum of the last three terms of \eqref{eq:involution-z-equation} minus $1$ be $V(t,z)$, and put
\[
 \Phi(t,u)=-\frac2\ell V(t,z_0+u).
\]
Then $z-z_0=\Phi(t,z-z_0)$, with positive $t$-order. Hence
\begin{equation}
 z_j(\ell)=\sum_{m=1}^{j}\frac1m[t^ju^{m-1}]\Phi(t,u)^m
 \quad(j\ge1).
 \label{eq:involution-z-finite}
\end{equation}
Each required $q_j$ is supplied by \eqref{eq:involution-B} and \eqref{eq:L}. Fixed-order substitution and the logarithmic derivative
$(\log A_+)'(x)=\tfrac12\log x+\bigO(x^{-1/2})$
prove the asymptotic remainder by Lemma~\ref{lem:residual}. In particular the coefficients and their required logarithmic derivatives stay bounded for large $\ell$, giving the stated $\bigO(X^{-1})$ remainder after $z_2$.
\end{proof}

\subsection{The exponentially small inverse displacement}

For the full interpolation \eqref{eq:involution-cont}, keep $x_+$ exact. A first sector correction is
\begin{equation}
 I^{-1}(y)=x_+-
 \frac{\cos(\pi x_+)A_-(x_+)/A_+(x_+)}{(\log A_+)'(x_+)}
        +\bigO(\ee^{-4\sqrt{x_+}}).
 \label{eq:involution-inverse-sector}
\end{equation}
The ratio of exact saddle functions satisfies
\[
 \frac{A_-(x)}{A_+(x)}
 \sim\ee^{-2\sqrt x}\frac{B(-x^{-1/2})}{B(x^{-1/2})}
 =\ee^{-2\sqrt x}\left(1-\frac7{12\sqrt x}+\cdots\right).
\]
Consequently,
\begin{align}
 I^{-1}(y)={}&x_+-\frac{2\cos(\pi x_+)}{\log x_+}\ee^{-2\sqrt{x_+}}
 \nonumber\\
 &+\bigO\!\left(\frac{\ee^{-2\sqrt{x_+}}}{\sqrt{x_+}\log x_+}\right)
 +\bigO(\ee^{-4\sqrt{x_+}}).
 \label{eq:involution-inverse-leading}
\end{align}
This follows from Taylor expansion or directly from Theorem~\ref{thm:transport}; all higher sectors follow from the same theorem. Products and derivatives generate higher harmonics in the phase $\pi x_+$ as well as powers of the stretched-exponential factor.

The distinction between $X$ and $x_+$ matters even inside the action. Since
$x_+-X\sim-2\sqrt X/\log X$,
\[
 \ee^{-2\sqrt{x_+}}
 =\ee^{-2\sqrt X}\exp\left(\frac2{\log X}
             +\bigO\!\left(\frac1{\sqrt X(\log X)^2}\right)\right).
\]
Thus replacing $x_+$ by $X$ without expanding the induced logarithmic correction discards an entire part of the sector amplitude. Replacing $\cos(\pi x_+)$ by $\cos(\pi X)$ is even more serious, because $x_+-X$ is unbounded. One must retain or expand the full phase, not just the exponential decay rate.

\section{Alternating permutations and prime-generated inverse sectors}
\label{sec:zigzag}

\subsection{Exact pole and Dirichlet decompositions}

Let $E_n$ count alternating permutations, with exponential generating function
\begin{equation}
 \sum_{n\ge0}E_n\frac{z^n}{n!}=\sec z+\tan z
   =\tan\left(\frac\pi4+\frac z2\right).
 \label{eq:zigzag-egf}
\end{equation}
This is A000111 \cite{oeiszigzag}. The equality to the shifted tangent makes its pole structure transparent: the poles are $\pi/2+2\pi k$, $k\in\mathbb Z$, and every residue is $-2$.

\begin{proposition}[Exact factorial--Dirichlet factorization]
\label{prop:zigzag-dirichlet}
Put
\begin{equation}
 A(x)=\frac4\pi\left(\frac2\pi\right)^x\Gamma(x+1).
 \label{eq:zigzag-A}
\end{equation}
For positive integers $n$,
\begin{equation}
 E_n=\begin{cases}
 A(n)\,\beta(n+1),&n\ \text{even},\\
 A(n)(1-2^{-n-1})\zeta(n+1),&n\ \text{odd},
 \end{cases}
 \label{eq:zigzag-parity}
\end{equation}
where $\beta(s)=\sum_{j\ge0}(-1)^j(2j+1)^{-s}$ is Dirichlet's beta function.
\end{proposition}
\begin{proof}
For $n\ge1$, integrate
$\tan(\pi/4+z/2)/z^{n+1}$ around large rectangles whose edges stay a fixed distance from the real poles. The tangent is uniformly bounded on these edges, while the denominator makes the contour integral tend to zero. The residue at zero is $E_n/n!$ and the other residues are $-2/(\pi/2+2\pi k)^{n+1}$. Absolute convergence of the resulting series gives
\[
 \frac{E_n}{n!}=2\sum_{k\in\mathbb Z}
                 \left(\frac\pi2+2\pi k\right)^{-n-1}.
\]
Factoring $(2/\pi)^{n+1}$ leaves the sum over all integers congruent to $1$ modulo $4$. Negative terms contribute a minus sign when $n$ is even, giving the alternating odd-integer sum $\beta(n+1)$. When $n$ is odd, they contribute positively and give the sum over all positive odd integers, $(1-2^{-n-1})\zeta(n+1)$. This proves the formula. The definition of the beta function agrees with \texttt{DirichletBeta} \cite{wlbeta}.
\end{proof}

Two useful smooth envelopes for $x>0$ are therefore
\begin{equation}
 E_{\mathrm e}(x)=A(x)\beta(x+1),\qquad
 E_{\mathrm o}(x)=A(x)(1-2^{-x-1})\zeta(x+1).
 \label{eq:zigzag-envelopes}
\end{equation}
They interpolate the even and odd subsequences respectively, not the full sequence. A full real interpolation is
\begin{equation}
 E(x)=A(x)\left\{
 \sum_{\substack{d\ge1\\d\equiv1\ (4)}}d^{-x-1}
 -\cos(\pi x)\sum_{\substack{d\ge1\\d\equiv3\ (4)}}d^{-x-1}
 \right\},\qquad x>0.
 \label{eq:zigzag-cont}
\end{equation}
The series converge absolutely. Since the bracket is $1+\bigO(3^{-x})$ with similarly small derivatives, each chosen function is eventually positive and increasing.

\subsection{The common factorial core and its inverse}

Set
\[
 a=\log\frac2\pi,\qquad C_A=\frac{4\sqrt{2\pi}}\pi.
\]
Stirling's expansion gives
\begin{equation}
 \log A(x)\sim x(\log x-1+a)+\frac12\log x+\log C_A
       +\sum_{m\ge1}\frac{B_{2m}}{2m(2m-1)x^{2m-1}}.
 \label{eq:zigzag-A-log}
\end{equation}
For $L=\log(y/C_A)$, define
\begin{equation}
 X=\frac{L}{W_0(2L/(\ee\pi))},\qquad
 \ell=\log X,\qquad s=\ell+a.
 \label{eq:zigzag-core}
\end{equation}
The positive increasing branch of $A$ then has the inverse expansion
\begin{align}
 A^{-1}(y)={}&X-\frac\ell{2s}
 +\frac1X\left(-\frac1{12s}+\frac\ell{4s^2}
                         -\frac{\ell^2}{8s^3}\right)
 +\bigO(X^{-2}).
 \label{eq:zigzag-A-inverse}
\end{align}
To prove it, substitute $x=X+d_0+d_1/X$ into \eqref{eq:zigzag-A-log}. The order-one equation gives $d_0=-\ell/(2s)$. At order $X^{-1}$, the terms are $sd_1+1/12+d_0/2+d_0^2/2$, giving the displayed $d_1$. Lemma~\ref{lem:residual} and the slope $s+o(1)$ justify the error. Arbitrarily high coefficients follow from the unbalanced gamma core and Proposition~\ref{prop:scaled}.

\subsection{Prime powers organize the logarithmic sectors}

For $\Re s>1$, Euler products give
\begin{align}
 \log\beta(s)&=\sum_{\substack{p\ \mathrm{prime}\\p\ne2}}
               \sum_{m\ge1}\frac{\chi_4(p)^m}{m p^{ms}},\nonumber\\
 \log((1-2^{-s})\zeta(s))&=
             \sum_{\substack{p\ \mathrm{prime}\\p\ne2}}
               \sum_{m\ge1}\frac1{m p^{ms}},
 \label{eq:zigzag-primes}
\end{align}
where $\chi_4(p)=1$ for $p\equiv1\pmod4$ and $-1$ for $p\equiv3\pmod4$. These identities can be proved directly by multiplying the absolutely convergent geometric factors over primes, using unique factorization, and expanding $-\log(1-z)=\sum_{m\ge1}z^m/m$.

Thus the primitive logarithmic sectors are indexed by odd prime powers, with actions $m\log p$. Exponentiating, differentiating, and reverting generates products indexed by odd composite integers. At any fixed action cutoff only finitely many such integers occur. This arithmetic structure is qualitatively different from the two-saddle involution case, but it is still locally finite.

Let $x_0=A^{-1}(y)$ be retained exactly and let
\[
 p_0=(\log A)'(x_0)=\psi(x_0+1)+\log(2/\pi).
\]
The first parity-dependent inverse layers are
\begin{align}
 E_{\mathrm e}^{-1}(y)&=x_0+\frac{3^{-x_0-1}}{p_0}
                  +\bigO\!\left(\frac{5^{-x_0}}{\log x_0}\right),\nonumber\\
 E_{\mathrm o}^{-1}(y)&=x_0-\frac{3^{-x_0-1}}{p_0}
                  +\bigO\!\left(\frac{5^{-x_0}}{\log x_0}\right).
 \label{eq:zigzag-inverse-parity}
\end{align}
For the full interpolation the first correction is instead
\begin{equation}
 E^{-1}(y)=x_0+\frac{\cos(\pi x_0)3^{-x_0-1}}{p_0}
            +\bigO\!\left(\frac{5^{-x_0}}{\log x_0}\right).
 \label{eq:zigzag-full-inverse}
\end{equation}
Here $5^{-x_0}$ dominates both the next integer sector and the square $9^{-x_0}$ of the first correction. These formulas follow from Theorem~\ref{thm:transport}, which also supplies every higher inverse coefficient. For the two envelopes, one may use the logarithmic sector sums \eqref{eq:zigzag-primes} directly as $h$ in that theorem.

The sector scale looks unusual when expressed entirely in $y$. Since $x_0\sim\log y/\log\log y$,
\begin{equation}
 d^{-x_0}=\exp\left[-(\log d)\frac{\log y}{\log\log y}(1+o(1))\right].
 \label{eq:zigzag-y-scale}
\end{equation}
For fixed $d>1$, this is smaller than every inverse power of $\log y$, but larger than every fixed positive power $y^{-\epsilon}$. It lies between familiar logarithmic and algebraic target scales. Keeping the core inverse in the action makes this hierarchy much easier to manipulate correctly.

\section{Connected labeled graphs: every exponential layer}
\label{sec:graphs}

\subsection{A formal generating function with zero analytic radius}

Let
\[
 a_n=2^{\binom n2}
\]
be the number of simple labeled graphs on $n$ vertices, and let $c_n$ be the number of connected such graphs. These are A006125 and A001187 \cite{oeisallgraphs,oeisconnected}. The exponential formula for connected components reads
\begin{equation}
 A(z)=\sum_{n\ge0}a_n\frac{z^n}{n!},\qquad
 \log A(z)=\sum_{n\ge1}c_n\frac{z^n}{n!}.
 \label{eq:graph-egf}
\end{equation}
These are formal power series. The first has radius of convergence zero, so analytic singularity transfer at a positive radius would be inapplicable. Nevertheless, fixed coefficient extraction and formal logarithms are entirely legitimate.

For completeness, a graph has a unique partition of its vertex labels into connected components. The exponential expansion of the formal series for connected graphs sums over such set partitions, with the factor $1/m!$ removing the ordering of $m$ components. This proves \eqref{eq:graph-egf} without making any analytic assertion.

Define the reciprocal coefficients
\begin{equation}
 b_k=k![z^k]\frac1{A(z)}.
 \label{eq:graph-b}
\end{equation}
They have the explicit finite partition formula
\begin{equation}
 b_k=k!\sum_{\substack{\sum_{j=1}^k jm_j=k\\M=\sum_{j=1}^k m_j}}
       (-1)^M M!\prod_{j=1}^k
            \frac1{m_j!}\left(\frac{2^{\binom j2}}{j!}\right)^{m_j},
 \qquad b_0=1.
 \label{eq:graph-b-finite}
\end{equation}
This follows by expanding $1/A=\sum_{M\ge0}(-1)^M(A-1)^M$. The first coefficients are
\begin{equation}
 b_0,b_1,b_2,b_3,b_4,b_5,b_6
   =1,-1,0,-2,-24,-544,-22320.
 \label{eq:graph-b-first}
\end{equation}

\subsection{A fixed-order exponential transseries theorem}

\begin{theorem}[Connected-graph expansion]
\label{thm:graphs}
For every fixed integer $K\ge0$,
\begin{equation}
 \frac{c_n}{2^{\binom n2}}
 =\sum_{k=0}^{K}\binom nk b_k
       2^{k(k+1)/2}\,2^{-kn}
  +\bigO_K\!\left(n^{K+1}2^{-(K+1)n}\right).
 \label{eq:graph-expansion}
\end{equation}
Thus every exponential coefficient is the polynomial
\begin{equation}
 P_k(x)=\binom xk b_k2^{k(k+1)/2}.
 \label{eq:graph-P}
\end{equation}
In particular,
\begin{equation}
 \frac{c_n}{2^{\binom n2}}
 =1-2n2^{-n}-128\binom n3 2^{-3n}
       -24576\binom n4 2^{-4n}
       +\bigO(n^5 2^{-5n}).
 \label{eq:graph-first}
\end{equation}
The entire $2^{-2n}$ sector vanishes.
\end{theorem}

\begin{proof}
Expand the formal logarithm in \eqref{eq:graph-egf} as
\[
 \log A(z)=\sum_{m\ge1}\frac{(-1)^{m-1}}m(A(z)-1)^m.
\]
A contribution to $[z^n]$ is an ordered composition of $n$ into $m$ positive block sizes. If one block has size $n-k>n/2$, it is unique. Its $m$ possible positions cancel the factor $1/m$. The other blocks have total size $k$ and their alternating sum is exactly $[z^k](1/A(z))$. Consequently the total contribution of all terms with a giant block of size $n-k$ is
\begin{equation}
 \binom nk b_k a_{n-k},\qquad 0\le k<n/2.
 \label{eq:graph-giant}
\end{equation}
Dividing by $a_n$ and using
\[
 \binom{n-k}{2}-\binom n2=-kn+\frac{k(k+1)}2
\]
gives the coefficients in \eqref{eq:graph-expansion}.

It remains to bound the omitted terms. For $k\ge1$, the absolute sum defining $b_k$ is bounded by
\[
 |b_k|\le k!\,2^{k-1}a_k.
\]
Indeed, there are $2^{k-1}$ ordered compositions of $k$; for each one the labeling factor is at most $k!$, and the product of internal graph counts is at most $a_k$. Hence the absolute relative contribution of \eqref{eq:graph-giant} is at most
\begin{equation}
 (2n)^k2^{-k(n-k)}.
 \label{eq:graph-tail-bound}
\end{equation}
For $K+1\le k\le n/4$, the ratio of successive bounds is
$2n\,2^{-n+2k+1}$, which is less than $1/2$ for large $n$. The sum is therefore bounded by a constant times its first term,
$\bigO_K(n^{K+1}2^{-(K+1)n})$. For $n/4<k<n/2$, the exponent $k(n-k)$ is at least $3n^2/16$ up to an inessential integer-endpoint error, so these terms are $2^{-c n^2+\bigO(n\log n)}$ and are negligible at every fixed exponential order.

Finally consider terms having no block larger than $n/2$. If their sizes are $j_1,\ldots,j_m$, then
\[
 \sum_i j_i^2\le\frac n2\sum_i j_i=\frac{n^2}2.
\]
The product of their graph counts divided by $a_n$ is at most $2^{-n^2/4}$. Bounding all labeling factors by $n!$ and the number of ordered compositions by $2^{n-1}$ bounds their total by
$2^{-n^2/4+\bigO(n\log n)}$. Combining the three bounds proves \eqref{eq:graph-expansion}. Substituting \eqref{eq:graph-b-first} proves \eqref{eq:graph-first}.
\end{proof}

This is an all-orders theorem in exponentially small lattice sectors. It is not a convergent-series assertion for a fixed real $x$. In particular, we have not defined a canonical connected-graph interpolation by summing the formal expression $\sum_kP_k(x)2^{-kx}$.

\subsection{A quadratic core and explicit inverse coefficients}

Let $a=\log2$. The logarithm of the all-graph count has the exact quadratic core
\begin{equation}
 H(x)=\frac a2x(x-1),\qquad
 X=\frac{1+\sqrt{1+8\log y/a}}2,
 \qquad p=H'(X)=a(X-\tfrac12).
 \label{eq:graph-core}
\end{equation}
There is no Lambert function at this stage. Define $t=2^{-X}$ and seek a formal range inverse
\begin{equation}
 x\sim X+\sum_{N\ge1}d_N(X)t^N.
 \label{eq:graph-inverse}
\end{equation}
All coefficients have a finite extraction formula. Put
\begin{equation}
 R_X(t,u)=\log\left(1+\sum_{k\ge1}
                 P_k(X+u)t^k\ee^{-kau}\right).
 \label{eq:graph-R}
\end{equation}
Then Proposition~\ref{prop:scaled}, in its additive form, gives
\begin{equation}
 d_N(X)=\sum_{m=1}^{N}\frac{(-1)^m}{m}
 [t^Nu^{m-1}]\left(\frac{R_X(t,u)}{p+au/2}\right)^m.
 \label{eq:graph-d}
\end{equation}
The denominator is simply $(H(X+u)-H(X))/u$. At exponential order $N$, only $P_1,\ldots,P_N$ are needed, each given explicitly by \eqref{eq:graph-b-finite} and \eqref{eq:graph-P}.

The first two coefficients are
\begin{align}
 d_1(X)&=\frac{2X}{p},\nonumber\\
 d_2(X)&=\frac{2X^2-2(aX-1)d_1(X)-(a/2)d_1(X)^2}{p}.
 \label{eq:graph-d-first}
\end{align}
To see the second formula directly, use
\[
 R_X(t,u)=-2Xt+2(aX-1)ut-2X^2t^2+\cdots
\]
in $pu+(a/2)u^2+R_X(t,u)=0$. Although $P_2=0$, the quadratic logarithm and the shift of the first sector generate a nonzero $d_2$. This is a concrete example of inverse sector creation.

\begin{corollary}[A rigorous inverse on the graph sequence's range]
\label{cor:graph-range}
At $y=c_n$, the truncation through order $N$ of \eqref{eq:graph-inverse} satisfies
\begin{equation}
 X+\sum_{j=1}^Nd_j(X)2^{-jX}
 =n+\bigO_N\!\left(n^N2^{-(N+1)n}\right).
 \label{eq:graph-range-error}
\end{equation}
No interpolation of $c_n$ is required for this statement.
\end{corollary}
\begin{proof}
For fixed $N$, use the smooth finite forward model
\[
 F_N(x)=\ee^{H(x)}\left(1+\sum_{k=1}^NP_k(x)2^{-kx}\right).
\]
It is positive and increasing for all sufficiently large $x$, with logarithmic derivative $ax+\bigO(1)$. Theorem~\ref{thm:graphs} gives a logarithmic discrepancy between $F_N(n)$ and $c_n$ of order $n^{N+1}2^{-(N+1)n}$. Lemma~\ref{lem:residual} therefore places the exact inverse $F_N^{-1}(c_n)$ within $\bigO_N(n^N2^{-(N+1)n})$ of $n$.

The formal coefficient cancellation in \eqref{eq:graph-d} gives the same error for its finite inverse truncation. One can bound the omitted terms directly: $P_k(x)=\bigO_k(x^k)$, $p\asymp X$, and \eqref{eq:graph-d} gives $d_j(X)=\bigO_j(X^{j-1})$. Taylor remainders for this finite model then have size $\bigO_N(X^{N+1}2^{-(N+1)X})$ in logarithmic phase, or $\bigO_N(X^N2^{-(N+1)X})$ after division by the slope. Since $X=n+\bigO(2^{-n})$, the same bound holds with $n$. Combining the estimates proves the result.
\end{proof}

Thus the inverse expansion is operationally meaningful even without a continuous connected-graph counting function. At sufficiently large exact counts, nearest-integer rounding recovers the vertex number. A threshold for an arbitrary target requires either bracketing consecutive exact counts or an interpolation together with the separation test from Section~\ref{sec:scope}.

\section{Necklaces and Lyndon words: arithmetic coefficients and accumulating actions}
\label{sec:necklaces}

\subsection{Exact divisor sums}

Fix an alphabet size $q\ge2$. The number of length-$n$ necklaces under cyclic rotation is
\begin{equation}
 N_q(n)=\frac1n\sum_{d\mid n}\varphi(d)q^{n/d},
 \label{eq:necklace}
\end{equation}
and the number of primitive necklaces, equivalently Lyndon words of length $n$ over the ordered alphabet, is
\begin{equation}
 L_q(n)=\frac1n\sum_{d\mid n}\mu(d)q^{n/d}.
 \label{eq:lyndon}
\end{equation}
The binary cases are A000031 and A001037 \cite{oeisnecklace,oeislyndon}.

For \eqref{eq:necklace}, Burnside averaging over the $n$ rotations counts $q^{\gcd(n,r)}$ words fixed by rotation through $r$ places. Grouping by $d=n/\gcd(n,r)$ gives $\varphi(d)$ rotations of the corresponding type. For \eqref{eq:lyndon}, let $P_q(n)$ be the number of aperiodic words. Every word has a unique least period dividing $n$, so $q^n=\sum_{d\mid n}P_q(d)$. M\"obius inversion gives $P_q(n)=\sum_{d\mid n}\mu(d)q^{n/d}$, and every aperiodic orbit has exactly $n$ words. This proves both formulas.

Factoring out $q^n/n$, write
\begin{equation}
 N_q(n)=\frac{q^n}{n}\left(1+
 \sum_{\substack{d\mid n\\d\ge2}}\varphi(d)\ee^{-\omega_dn}\right),
 \qquad \omega_d=(\log q)\left(1-\frac1d\right).
 \label{eq:necklace-sectors}
\end{equation}
For Lyndon words, replace $\varphi(d)$ by $\mu(d)$. Each divisor contributes a different exponential action.

\begin{proposition}[Uniform finite action cutoffs]
\label{prop:necklace-cutoff}
For every fixed integer $D\ge2$,
\begin{align}
 \frac{nN_q(n)}{q^n}
 &=1+\sum_{d=2}^{D}\varphi(d)\,\mathbf1_{d\mid n}
             q^{-n(1-1/d)}
    +\bigO\!\left(nq^{-n(1-1/(D+1))}\right),\nonumber\\
 \frac{nL_q(n)}{q^n}
 &=1+\sum_{d=2}^{D}\mu(d)\,\mathbf1_{d\mid n}
             q^{-n(1-1/d)}
    +\bigO\!\left(nq^{-n(1-1/(D+1))}\right).
 \label{eq:necklace-cutoff}
\end{align}
\end{proposition}
\begin{proof}
Every omitted divisor satisfies $d\ge D+1$, hence
$1-1/d\ge1-1/(D+1)$. In the necklace sum use
$\sum_{d\mid n}\varphi(d)=n$. In the primitive sum use
$\sum_{d\mid n}|\mu(d)|\le\sum_{d\mid n}1\le n$.
The stated bounds follow immediately.
\end{proof}

\subsection{Why constant-coefficient exponentials are insufficient}

Let $\epsilon_n=nN_q(n)/q^n-1$. The first candidate action is $(\log q)/2$. But Proposition~\ref{prop:necklace-cutoff} gives
\begin{equation}
 q^{n/2}\epsilon_n\longrightarrow
 \begin{cases}1,&n\to\infty\ \text{through even integers},\\
               0,&n\to\infty\ \text{through odd integers}.
 \end{cases}
 \label{eq:necklace-parity-obstruction}
\end{equation}
For primitive necklaces the even limit is $-1$. Thus no single constant coefficient can describe the first exponential correction on all integers. Arithmetic indicators are essential, not negligible nuisances.

Moreover,
\begin{equation}
 \omega_2<\omega_3<\cdots<\log q,
 \qquad \omega_d\longrightarrow\log q.
 \label{eq:necklace-action-accumulation}
\end{equation}
There are infinitely many primitive actions below the finite action $\log q$. In particular, the nonlinear square of the first sector has action $2\omega_2=\log q$, but infinitely many elementary divisor sectors precede it. A finite grid of positive generators cannot have this property: its support below every fixed bound is finite.

This is an obstruction to the \emph{locally finite finite-grid calculus}, not a proof that no generalized transseries can exist. A well-ordered support can have order type $\omega$ and accumulate at a cut; larger Hahn-type or arithmetic-coefficient constructions may accommodate it. What fails is the convenient claim that finitely many sector coefficients suffice to reach every bounded action.

\subsection{Explicit finite-cutoff interpolations and inverses}

The dominant counting model is $f(x)=q^x/x$. Put $a=\log q$. Its large inverse is
\begin{equation}
 X=-\frac1a W_{-1}\!\left(-\frac ay\right),
 \qquad aX-\log X=\log y.
 \label{eq:necklace-core}
\end{equation}
This core is exact for the dominant model and again requires the negative Lambert branch.

To give a precise meaning to finite arithmetic corrections off the integers, define the real trigonometric filter
\begin{equation}
 D_d(x)=\frac1d\sum_{r=0}^{d-1}\cos\left(\frac{2\pi r x}{d}\right).
 \label{eq:divisibility-filter}
\end{equation}
At integer $x=n$, the finite geometric sum proves $D_d(n)=\mathbf1_{d\mid n}$. For fixed $D$, set
\begin{equation}
 F_{q,D}^{(w)}(x)=\frac{q^x}{x}
       \left(1+\sum_{d=2}^{D}w_dD_d(x)\ee^{-\omega_dx}\right),
 \label{eq:necklace-finite-model}
\end{equation}
where $w_d=\varphi(d)$ or $w_d=\mu(d)$. This is an explicit, eventually positive and increasing function: the finite correction and its derivative tend exponentially to zero, while $(\log f)'=a-1/x$ tends to $a>0$.

Every inverse sector of this \emph{fixed finite model} follows from Theorem~\ref{thm:transport}, with $g$ the exact Lambert inverse of the dominant phase. Its first displacement is
\begin{equation}
 (F_{q,D}^{(w)})^{-1}(y)
 =X-\frac{w_2D_2(X)}{a-1/X}q^{-X/2}
          +\bigO(q^{-2X/3}),
 \label{eq:necklace-finite-inverse}
\end{equation}
where $D_2(X)=(1+\cos\pi X)/2$. For $D=2$, the actual next nonlinear term is smaller, of order $q^{-X}$; the displayed weaker bound remains valid. For $D\ge3$, a divisor-$3$ sector may occur at order $q^{-2X/3}$.

At the true sequence values, Proposition~\ref{prop:necklace-cutoff} and Lemma~\ref{lem:residual} imply the rigorous range approximations
\begin{align}
 (F_{q,D}^{(\varphi)})^{-1}(N_q(n))
 &=n+\bigO\!\left(nq^{-n(1-1/(D+1))}\right),\nonumber\\
 (F_{q,D}^{(\mu)})^{-1}(L_q(n))
 &=n+\bigO\!\left(nq^{-n(1-1/(D+1))}\right).
 \label{eq:necklace-range}
\end{align}
Indeed, the relative forward error has exactly that size, and the logarithmic derivative of the finite model stays bounded away from zero. These estimates justify finite-cutoff inverse computations without inventing an infinite real interpolation.

It would be unjustified to replace $D$ by infinity in \eqref{eq:necklace-finite-model} without a convergence analysis: for real noninteger $x$, the filters no longer restrict the sum to divisors of a single integer. The finite divisor identity has become an unrelated infinite series. This is a concrete place where a formal manipulation of a combinatorial formula can silently change the problem.

\section{Harmonic functions: exponential and endpoint-Puiseux inversion}
\label{sec:harmonic}

\subsection{A slowly increasing combinatorial function}

Harmonic numbers occur throughout permutation enumeration, including the $k=2$ fixed-cycle Stirling formula. Their standard real continuation is
\begin{equation}
 H(x)=\psi(x+1)+\gamma,\qquad x\ge0,
 \label{eq:harmonic-cont}
\end{equation}
which agrees with $H_n=\sum_{j=1}^n1/j$ at integers. The digamma and generalized-harmonic conventions agree with \texttt{HarmonicNumber} \cite{wlharmonic}. Since
$H'(x)=\sum_{j\ge0}(x+1+j)^{-2}>0$, this function has a unique increasing inverse on $[0,\infty)$.

The centered variable $w=x+1/2$ eliminates odd inverse powers from the logarithmic correction. The shifted digamma expansion gives
\begin{equation}
 H(x)-\gamma\sim\log w+C(w^{-2}),\qquad
 C(z)=-\sum_{m\ge1}\frac{B_{2m}(1/2)}{2m}z^m.
 \label{eq:harmonic-centered}
\end{equation}
For example, $C(z)=z/24-7z^2/960+\cdots$. This follows by differentiating the shifted gamma expansion, or by Euler--Maclaurin summation with the half-integer shift \cite{dlmfgamma,dlmfbernoulli}.

\begin{theorem}[All coefficients of the harmonic inverse]
\label{thm:harmonic-inverse}
Let $X=\ee^{y-\gamma}$. Then
\begin{equation}
 H^{-1}(y)\sim X-\frac12+\sum_{m\ge1}h_mX^{1-2m},
 \qquad
 h_m=-\frac1{2m-1}[z^m]\exp\bigl((2m-1)C(z)\bigr).
 \label{eq:harmonic-inverse}
\end{equation}
The first five coefficients are
\begin{align}
 h_1&=-\frac1{24},&h_2&=\frac3{640},&
 h_3&=-\frac{1525}{580608},\nonumber\\
 h_4&=\frac{615881}{199065600},&
 h_5&=-\frac{3058641}{504627200}.
 \label{eq:harmonic-coeff}
\end{align}
A truncation through $h_MX^{1-2M}$ has error $\bigO(X^{-2M-1})$.
\end{theorem}
\begin{proof}
Put $v=1/w$ and $t=1/X$. Equation~\eqref{eq:harmonic-centered} becomes formally
\[
 v=t\exp(C(v^2)).
\]
The Laurent form of Lagrange inversion gives, for $m\ge1$,
\[
 [t^{2m-1}]\frac1{v(t)}
 =-\frac1{2m-1}[u^{2m}]
             \exp\bigl((2m-1)C(u^2)\bigr),
\]
which is the coefficient in \eqref{eq:harmonic-inverse}. One may derive this Laurent version by applying the usual residue proof of Lagrange inversion to $v^{-1}$; the leading $t^{-1}$ term is separated before extracting positive powers. Oddness of $v(t)$ forces the displayed parity of powers in $w=1/v$.

Substitution of a finite truncation into a sufficiently long centered digamma expansion gives a forward residual of order $X^{-2M-2}$. Since $H'(x)\asymp X^{-1}$ near the proposed inverse, Lemma~\ref{lem:residual} gives an inverse error $\bigO(X^{-2M-1})$. Formula~\eqref{eq:E} computes the displayed rational coefficients exactly.
\end{proof}

In the target variable, the correction sequence is
\[
 \ee^{y-\gamma},\quad 1,\quad
 \ee^{-(y-\gamma)},\quad\ee^{-3(y-\gamma)},\quad\ldots.
\]
Thus even an inverse-power expansion in the centered forward variable becomes an exponential inverse series. The Bernoulli coefficient series remains generally divergent; no convergence claim is made by writing all its coefficients.

\subsection{Generalized harmonic functions near their finite limit}

Fix a real $s>1$ and define
\begin{equation}
 H_s(x)=\zeta(s)-\zeta(s,x+1),\qquad x\ge0.
 \label{eq:general-harmonic}
\end{equation}
The Hurwitz-zeta series shows
$H_s'(x)=s\zeta(s+1,x+1)>0$. Its range approaches the finite endpoint $\zeta(s)$, so the inverse asymptotic problem is now $y\uparrow\zeta(s)$, not $y\to\infty$.

Put
\begin{equation}
 p=s-1,\qquad \delta=\zeta(s)-y,\qquad
 T=(p\delta)^{-1/p},\qquad w=x+\frac12.
 \label{eq:general-harmonic-core}
\end{equation}
The centered Hurwitz-zeta expansion \cite{dlmfhurwitz} gives
\begin{align}
 \zeta(s,x+1)&\sim\frac{w^{-p}}pD_s(w^{-2}),\nonumber\\
 D_s(z)&=1+\sum_{m\ge1}
  \frac{pB_{2m}(1/2)\rising{s}{2m-1}}{(2m)!}z^m.
 \label{eq:general-harmonic-D}
\end{align}
In particular $D_s(z)=1-ps\,z/24+\cdots$.

\begin{theorem}[Endpoint inverse coefficients]
\label{thm:general-harmonic-inverse}
For fixed $s>1$,
\begin{equation}
 H_s^{-1}(y)\sim T-\frac12+\sum_{m\ge1}h_{s,m}T^{1-2m},
 \quad
 h_{s,m}=-\frac1{2m-1}[z^m]
                    D_s(z)^{-(2m-1)/p}.
 \label{eq:general-harmonic-inverse}
\end{equation}
In particular,
\begin{equation}
 H_s^{-1}(y)=T-\frac12-\frac{s}{24T}+\bigO(T^{-3}).
 \label{eq:general-harmonic-first}
\end{equation}
The remainder after order $M$ is $\bigO(T^{-2M-1})$.
\end{theorem}
\begin{proof}
The forward equation reads $T=wD_s(w^{-2})^{-1/p}$. Thus with $v=1/w$ and $t=1/T$,
\[
 v=tD_s(v^2)^{-1/p}.
\]
The same Laurent Lagrange formula used for the harmonic inverse gives \eqref{eq:general-harmonic-inverse}. The first coefficient is $-s/24$ by direct expansion. For a finite truncation, the relative forward error in $\delta$ is $\bigO(T^{-2M-2})$, hence the absolute error is $\bigO(T^{-p-2M-2})$. The derivative magnitude of the tail is $\asymp T^{-p-1}$, so the inverse error is $\bigO(T^{-2M-1})$.
\end{proof}

Because $T$ is a fractional power of the distance $\delta$ to the endpoint, this is a Puiseux-type expansion in $\delta^{1/p}$. When $p$ is not rational, the powers need not be rational; ``generalized Puiseux'' is the appropriate description. The finite coefficient formula still makes sense for every fixed real $s>1$.

\section{Synthesis: scales, closure, and what has been proved}
\label{sec:synthesis}

\subsection{A comparative map}

Table~\ref{tab:classes} summarizes the different dominant balances. A symbol such as $X$ in an exponential action always denotes the appropriate exact or formally specified dominant inverse, not an interchangeable rough approximation.

\begin{table}[htbp]
\centering
\small
\renewcommand{\arraystretch}{1.25}
\begin{tabular}{P{34mm}P{46mm}P{64mm}}
\toprule
Family & Forward structure & Inverse structure\\\midrule
Catalan, Fuss--Catalan, fixed-height rectangles
 & $C\rho^x x^\beta\sum a_jx^{-j}$
 & $W_{-1}$ core and inverse powers of that core\\
Fixed-column Stirling II
 & Finite sum of real exponentials
 & Convergent multivariate generalized powers of $1/y$\\
Fixed-column Stirling I
 & $\Gamma(x)$ times polynomials in $\log x$
 & $W_0$ core; $\log\log X/\log X$ and rational-logarithmic corrections\\
Motzkin
 & Two exact endpoint moments; relative $3^{-x}\cos\pi x$
 & $W_{-1}$ core; exact dominant inverse with oscillatory exponential corrections\\
Involutions
 & $(x/\ee)^{x/2}\ee^{\sqrt x}$; relative $\ee^{-2\sqrt x}\cos\pi x$
 & $W_0$ core; half-powers with rational-log coefficients, then stretched-exponential sectors\\
Alternating permutations
 & $\Gamma(x+1)(2/\pi)^x$ times odd Dirichlet sums
 & Factorial core and prime-generated exponential sectors\\
Connected graphs
 & $2^{x(x-1)/2}\sum P_k(x)2^{-kx}$ on integers
 & Quadratic core; rational-function sector coefficients; rigorous range inverse\\
Necklaces, Lyndon words
 & $q^n/n$ times divisor-dependent sectors
 & $W_{-1}$ core; arithmetic finite-cutoff inverses with accumulating actions\\
Harmonic functions
 & Logarithmic growth, or convergence to $\zeta(s)$
 & Exponential inverse, or endpoint generalized Puiseux expansion\\
\bottomrule
\end{tabular}
\caption{Dominant structures and their inverse scales. Delannoy and Schr\"oder inverses fall in the first row's balanced class, with the moment interpolation specified in Section~\ref{sec:moments}.}
\label{tab:classes}
\end{table}

The negative Lambert branch is not a rare special case. Any balanced exponentially growing sequence with an algebraic prefactor $x^\beta$, $\beta<0$, leads to it. By contrast, factorial growth has a positive Lambert core because the $x\log x$ term changes the dominant balance. Connected graph enumeration is faster still in logarithmic phase, but its inverse core is simpler: a quadratic root.

\subsection{Why nonlinear inversion creates sectors but not arbitrary constants}

\begin{proposition}[Additive support closure under inverse transport]
\label{prop:support}
Suppose an exact small correction, expressed at the exact dominant inverse $x_0$, has sectors
\[
 \varepsilon(x_0)=\sum_{j=1}^r
       a_j(x_0)\exp(-\lambda_j\Omega(x_0)),\qquad\lambda_j>0.
\]
Assume the coefficient class is closed under products and under
$\mathcal D=((\log F_0)'(x_0))^{-1}\dd/\dd x_0$, including the factors introduced by differentiating $\Omega$. Then the formal inverse produced by Theorem~\ref{thm:transport} has exponential support contained in
\[
 \left\{\sum_{j=1}^rm_j\lambda_j:
              m_j\in\mathbb Z_{\ge0},\quad\sum_jm_j>0\right\}.
\]
At every finite action cutoff, the coefficients are finite sums. The same conclusion holds for any locally finite positive action set.
\end{proposition}
\begin{proof}
Expanding $\log(1+\varepsilon)$ produces only products of input sectors, so its actions are sums of the $\lambda_j$. Differentiation multiplies a sector by coefficient factors and does not change its exponential rate in the stated coefficient class. Formula~\eqref{eq:transport} uses only these two operations. Positive local finiteness makes extraction below a fixed action cutoff finite. No new independent constant appears in any operation.
\end{proof}

The proposition explains the nonzero second inverse graph sector, the composite odd-integer sectors of the alternating-permutation inverse, and the higher involution actions $2m\sqrt{x_+}$. It does not manufacture Stokes constants for a dominant asymptotic expansion that never specified a small sector in the first place. Those constants belong to the choice and analysis of the original function, not to formal reversion alone.

For oscillatory coefficients, multiplication adds frequencies and differentiation retains them. Thus nonlinear inversion can generate $\cos(2\pi x_0)$ and $\sin(2\pi x_0)$ even if the forward interpolation had only $\cos(\pi x)$. Treating an oscillatory prefactor as a constant during differentiation gives incorrect second and higher layers.

\subsection{Status of the results}

The following distinctions summarize the logical scope of the article.

\paragraph{Convergent identities.}
The finite-exponential inverse of Theorem~\ref{thm:finite-exp} is locally convergent in its small variables. The Motzkin beta representations converge absolutely. The alternating-permutation Dirichlet sums and Euler products converge absolutely in the stated half-plane. The finite enumeration and divisor formulas are exact.

\paragraph{All-orders asymptotic expansions with error control.}
The gamma-quotient, fixed-column Stirling I, endpoint-moment, moving-saddle, and harmonic results have a remainder at every prescribed finite algebraic order. Their inverse errors follow from a quantified slope and a forward residual. These statements do not assert convergence of the corresponding infinite Bernoulli or saddle series.

\paragraph{Exactly normalized exponentially small sectors.}
Motzkin and involution sectors are defined by exact separate integrals. Alternating-permutation sectors are defined by exact Dirichlet sums. Their inverse transport formulas refer to the exact dominant inverse. This gives a meaningful all-sector formal expansion without confusing the error of a fixed algebraic truncation with the size of a subdominant sector.

\paragraph{Lattice transseries and range inverses.}
Theorem~\ref{thm:graphs} proves every fixed exponential layer for connected graph counts, and Corollary~\ref{cor:graph-range} proves the corresponding inverse on the range. The necklace results prove finite-cutoff arithmetic expansions and finite-model range inverses. Neither argument asserts a unique global real interpolation of the sequence.

\paragraph{Not claimed.}
No global resurgence or Borel summation theorem is established for all the examples. No uniform expansion is asserted as a fixed combinatorial parameter grows with $n$. No lower-endpoint Stokes coefficient is supplied for the Delannoy or Schr\"oder moment functions. No infinite noninteger divisor interpolation is assumed for necklaces. These are separate questions rather than consequences of the formulas already proved.

\subsection{A reproducible route from an enumeration formula to its inverse}

The derivations suggest a practical workflow. First choose an exact representation: a gamma quotient, finite exponential sum, positive moment, formal component generating function, or divisor sum. Then decide whether the desired output is a continuous inverse, a range inverse, or a staircase threshold. This prevents interpolation conventions from being inserted accidentally in the middle of an asymptotic calculation.

Next extract the logarithmic dominant phase and invert it as accurately as possible in closed form. The useful cores here are $ax+\beta\log x$, $\kappa x\log x+\lambda x$, and a quadratic polynomial. Keeping the core exact avoids repeated expansion of its internal logarithms. Normalize the remaining correction so that it has positive order, using a relative displacement if the absolute correction grows.

Then compute only the coefficients required by the chosen precision. Formula~\eqref{eq:E} handles exponentiation, \eqref{eq:L} handles logarithms, and the relevant finite Lagrange formula handles reversion. For multiple sectors, specify an action cutoff before truncating algebraic amplitudes. For arithmetic actions with a finite accumulation point, specify a finite divisor cutoff instead of claiming a complete bounded-action truncation.

Finally substitute the inverse approximation into the \emph{forward} equation. The residual is generally easier to estimate than the inverse error directly. Divide it by a justified slope bound, and use a rounding interval rather than an unqualified ceiling. This last step turns formal coefficient algebra into a trustworthy numerical or discrete conclusion.

\subsection{Further research directions}

Several concrete extensions are suggested by the results, but are not used as premises above.

One direction is to obtain exponentially improved remainders for the exact positive moment functions. For Motzkin and involution numbers, the small sector is already normalized by a real integral, so a useful next theorem would compare optimal truncation of the dominant amplitude directly with the known second sector. For Delannoy and Schr\"oder functions, the first task is instead to determine a contour- or summation-specific subdominant decomposition before proposing an inverse Stokes analysis.

A second direction is a uniform two-parameter theory. The fixed-$k$ Stirling expansions change when $k$ grows, and fixed-height rectangular tableaux change when the height is comparable with the width. The coefficient engines remain formally available, but the present remainder arguments are not uniform in those regimes. Identifying the transition scales and the appropriate dominant inverse would yield genuinely different asymptotic problems.

A third direction is an arithmetic transseries algebra for divisor sums. It should retain divisibility indicators, permit the action accumulation at $\log q$, and state exactly which products and inverse operations remain summable. The finite-cutoff estimates in Proposition~\ref{prop:necklace-cutoff} provide a concrete benchmark for such a theory. Merely replacing indicators by analytic expressions and summing indefinitely is not an adequate definition.

Finally, the connected-component proof suggests a general theorem for labeled structures whose total counts have a strongly convex superexponential logarithm. A unique giant component can again dominate formal logarithmic extraction. The appropriate hypotheses should control both the smaller-component reciprocal coefficients and the contributions of partitions with no giant block. Under such hypotheses one can seek analogues of \eqref{eq:graph-expansion} and then apply the same inverse coefficient transport. The graph theorem supplies a complete worked case, not a proof for every superexponential labeled class.

\appendix
\section{Coefficient formulas collected in one place}
\label{app:coefficients}

The formulas below show explicitly that the constructions are not dependent on unlisted auxiliary recurrences. Recurrences may be more efficient in a computer implementation, but a finite coefficient or partition formula determines every requested order.

\begin{longtable}{P{43mm}P{104mm}}
\toprule
Operation or family & Finite coefficient specification\\\midrule
\endfirsthead
\toprule
Operation or family & Finite coefficient specification\\\midrule
\endhead
Exponential amplitude &
$\displaystyle [t^n]\ee^{\sum q_jt^j}=\cE_n(q)$, with the partition sum \eqref{eq:E}.\\[3pt]
Logarithmic amplitude &
$\displaystyle [t^n]\log(1+\sum a_jt^j)=\cL_n(a)$, with \eqref{eq:L}.\\[3pt]
Gamma quotient &
Bernoulli-polynomial sum \eqref{eq:gamma-q}; normalized amplitudes $\cE_n(q)$.\\[3pt]
Balanced inverse &
$\displaystyle d_n=\sum_{m=1}^n\frac1m[t^nu^{m-1}]\Psi(t,u)^m$, equation~\eqref{eq:balanced-d}.\\[3pt]
Finite exponential inverse &
Multi-index generalized-binomial sum \eqref{eq:finite-exp-inverse}.\\[3pt]
Fixed-column Stirling I &
$\displaystyle P_{k,m}(\ell)=[z^{k-1}]\ee^{\ell z}\cE_m(r(z))/\Gamma(1+z)$; inverse by scaled Lagrange extraction.\\[3pt]
Moment endpoint &
$\displaystyle a_m=\sum_{j=0}^m\binom\theta j(-\kappa)^j(\nu)_jR_{m-j}(\alpha,\alpha+\nu+j)$.\\[3pt]
Involution saddle &
Gaussian partition sum \eqref{eq:involution-S}, with prefactor \eqref{eq:involution-B}. Only weights $\sum(k-2)m_k\le N$ contribute through $t^N$.\\[3pt]
Involution inverse &
After removing $z_0=-2/\ell$, finite extraction \eqref{eq:involution-z-finite}.\\[3pt]
Exact sector transport &
$\displaystyle \sum_{m\ge1}\frac{(-1)^m}{m!}\partial_L^{m-1}(g'h^m)$, truncated at the specified action.\\[3pt]
Connected graphs &
Reciprocal partition sum \eqref{eq:graph-b-finite}; polynomial $P_k$ in \eqref{eq:graph-P}; inverse extraction \eqref{eq:graph-d}.\\[3pt]
Harmonic inverse &
$\displaystyle h_m=-\frac{[z^m]\ee^{(2m-1)C(z)}}{2m-1}$.\\[3pt]
Generalized harmonic inverse &
$\displaystyle h_{s,m}=-\frac{[z^m]D_s(z)^{-(2m-1)/(s-1)}}{2m-1}$.\\
\bottomrule
\end{longtable}

For the moment sequences, a useful compact parameter dictionary for \eqref{eq:endpoint-coeff} is
\begin{center}
\begin{tabular}{lcccc}
\toprule
Sector & $\alpha$ & $\nu$ & $\theta$ & $\kappa$\\\midrule
Motzkin positive & $1$ & $3/2$ & $1/2$ & $3/4$\\
Motzkin negative & $1$ & $3/2$ & $1/2$ & $1/4$\\
Delannoy & $1$ & $1/2$ & $-1/2$ & $(3+2\sqrt2)/(4\sqrt2)$\\
Large Schr\"oder & $0$ & $3/2$ & $1/2$ & $(3+2\sqrt2)/(4\sqrt2)$\\\bottomrule
\end{tabular}
\end{center}
The normalization constants are outside this dictionary and are specified in the relevant section. Mixing those constants with normalized amplitude coefficients is a common source of inverse-core errors.

\section{Numerical and symbolic verification}
\label{app:verification}

The accompanying \texttt{verify.py} uses exact SymPy rational arithmetic for the coefficient identities and \texttt{mpmath} with 100 decimal digits for numerical evaluations. The companion \texttt{audit\_symbolic.py} checks several formulas independently. Both scripts are included with their actual output reports. These finite checks complement, rather than replace, the proofs.

\Needspace{12\baselineskip}
\subsection{Catalan inverse}

For each listed integer $n$, the target is the exact Catalan number $y=C_n$. The first column compares the Lambert core \eqref{eq:catalan-core} with $n$; the second uses all three corrections in \eqref{eq:catalan-inverse}.
\begin{center}
\begin{tabular}{rcc}
\toprule
$n$ & $|X-n|$ & Absolute error after three corrections\\\midrule
20 & $4.1976\times10^{-2}$ & $1.5033\times10^{-6}$\\
100 & $8.1677\times10^{-3}$ & $2.2520\times10^{-9}$\\
1000 & $8.1203\times10^{-4}$ & $2.2216\times10^{-13}$\\\bottomrule
\end{tabular}
\end{center}
The observed improvement is consistent with the proved $\bigO(X^{-4})$ inverse remainder. The exact gamma interpolation is specified; the experiment is not testing an arbitrary interpolation of the same integers.

\subsection{Involution forward and inverse checks}

The forward approximation uses the positive-saddle leading factor and the six amplitude coefficients through $x^{-5/2}$ in \eqref{eq:involution-B-first}. It is compared with the exact integer sum \eqref{eq:involution-enumeration}; hence its error includes the omitted negative saddle as well as the algebraic truncation. The inverse uses $z_0,z_1,z_2$ in \eqref{eq:involution-z} at $y=I_n$.
\begin{center}
\begin{tabular}{rcc}
\toprule
$n$ & Relative forward error & Absolute inverse error\\\midrule
100 & $1.3669\times10^{-8}$ & $4.9336\times10^{-4}$\\
400 & $1.6562\times10^{-10}$ & $1.0608\times10^{-4}$\\
1600 & $2.4275\times10^{-12}$ & $2.2684\times10^{-5}$\\\bottomrule
\end{tabular}
\end{center}
The independent symbolic audit substitutes the displayed inverse into \eqref{eq:involution-z-equation} and verifies exact cancellation through $t^2$. Thus the half-power coefficients are checked algebraically, not merely fitted to numerical values.

\subsection{The regenerated graph sector}

The connected graph counts are computed independently by the exact component-of-vertex-$1$ recurrence
\begin{equation}
 c_n=2^{\binom n2}-\sum_{j=1}^{n-1}
         \binom{n-1}{j-1}c_j2^{\binom{n-j}{2}}.
 \label{eq:graph-check-recurrence}
\end{equation}
This follows by choosing the vertex set of the connected component containing label $1$, selecting its connected graph, and placing an arbitrary graph on the remaining labels. It is used only as a verification algorithm; the article's asymptotic coefficients have the nonrecursive formula \eqref{eq:graph-b-finite}.

At $y=c_n$, the successive inverse errors are
\begin{center}
\begin{tabular}{rccc}
\toprule
$n$ & Quadratic core & One sector & Two sectors\\\midrule
12 & $7.3731\times10^{-4}$ & $1.8669\times10^{-6}$ & $6.4675\times10^{-8}$\\
20 & $2.8223\times10^{-6}$ & $4.8514\times10^{-11}$ & $1.0455\times10^{-14}$\\
40 & $2.6575\times10^{-12}$ & $9.1870\times10^{-23}$ & $3.8930\times10^{-32}$\\\bottomrule
\end{tabular}
\end{center}
The dramatic second improvement illustrates why a vanished forward sector must not simply be deleted from the inverse calculation.

The audit also evaluates graph forward remainders as exact rational numbers, avoiding cancellation at insufficient numerical precision. For example, after retaining sectors through $K=2$, the ratio of the remainder to $n^3 2^{-3n}$ is approximately $-18.2542$, $-19.7600$, and $-20.5400$ at $n=20,40,80$. These values approach the predicted leading coefficient $-128/6=-64/3$.

\subsection{A convergent inverse and an exponential inverse}

For $S_3^{-1}$, the complete multi-index formula is truncated at total degree six. For $H^{-1}$, the terms through $X^{-5}$ are used.
\begin{center}
\begin{tabular}{rcc}
\toprule
$n$ & Stirling-II logarithmic core error & Degree-six inverse error\\\midrule
10 & $4.8582\times10^{-2}$ & $9.6875\times10^{-13}$\\
20 & $8.2157\times10^{-4}$ & $3.9637\times10^{-25}$\\
50 & $4.2827\times10^{-9}$ & $4.1486\times10^{-62}$\\\bottomrule
\end{tabular}
\end{center}
\begin{center}
\begin{tabular}{rc}
\toprule
$n$ & Harmonic inverse absolute error at $y=H_n$\\\midrule
10 & $2.1550\times10^{-10}$\\
100 & $2.9870\times10^{-17}$\\
1000 & $3.0830\times10^{-24}$\\\bottomrule
\end{tabular}
\end{center}
The first experiment uses a genuinely convergent inverse series. The second uses a finite truncation of a generally divergent asymptotic series. Similar numerical success does not erase that analytic distinction.

\subsection{Additional independent checks and reproducibility}

The symbolic audit verifies the first two Motzkin logarithmic coefficients, the Stirling gamma-polynomial identity through $n=8$, and the generalized-harmonic inverse coefficient $-s/24$. Direct high-precision integration checks the Motzkin, Delannoy, Schr\"oder, and two-Gaussian involution moment identities for $n=0,\ldots,5$, with observed absolute discrepancies below $10^{-85}$. These are independent checks of the exact moment normalizations, which strongly affect both leading constants and inverse cores.

To reproduce the calculations, use Python 3.10 or later with SymPy and mpmath installed, then run
\begin{lstlisting}
python verify.py
python audit_symbolic.py
\end{lstlisting}
The scripts write \nolinkurl{verification_results.json}, \nolinkurl{verification_report.txt}, and \nolinkurl{symbolic_audit_report.txt} in their own directory. Exact assertions raise an error on failure. No internet access, OEIS download, or proprietary software is needed to run them.

\clearpage
\section{Wolfram Language coefficient recipes}
\label{app:wolfram}

The user-facing Wolfram functions \texttt{CatalanNumber}, \texttt{StirlingS1}, \texttt{StirlingS2}, \texttt{HarmonicNumber}, \texttt{DirichletBeta}, and \texttt{ProductLog} provide useful reference definitions \cite{wlcatalan,wlstirling1,wlstirling2,wlharmonic,wlbeta,wlproductlog}. The following symbolic recipes implement the coefficient formulas rather than asking a numerical inverse solver to infer an asymptotic expansion. They are supplied as a translation of the proved formulas; the executed checks accompanying this article used Python, not a Wolfram kernel.

Each row of \texttt{data} below is $\{a_i,b_i,\epsilon_i\}$ for a gamma factor $\Gamma(a_ix+b_i)^{\epsilon_i}$.
\begin{lstlisting}
ClearAll[gammaLogCoefficient, balancedInverseCoefficient,
  harmonicInverseCoefficient, catalanCore];

gammaLogCoefficient[data_List, r_Integer] /; r >= 1 :=
  (-1)^(r + 1)/(r (r + 1)) Total[
    (#[[3]] BernoulliB[r + 1, #[[2]]]/#[[1]]^r) & /@ data
  ];

balancedInverseCoefficient[q_List, a_, beta_, n_Integer] /;
    n >= 1 && Length[q] >= n :=
  Module[{t, u, psi},
    psi = -(beta Log[1 + t u] +
       Sum[q[[j]] (t/(1 + t u))^j, {j, 1, n}])/a;
    Simplify[Total[Table[
      SeriesCoefficient[
        SeriesCoefficient[psi^m, {t, 0, n}],
        {u, 0, m - 1}]/m,
      {m, 1, n}]]]
  ];

harmonicInverseCoefficient[m_Integer] /; m >= 1 :=
  Module[{z, c},
    c = -Sum[BernoulliB[2 j, 1/2] z^j/(2 j), {j, 1, m}];
    -SeriesCoefficient[Exp[(2 m - 1) c], {z, 0, m}]/(2 m - 1)
  ];

catalanCore[y_] := -3/(2 Log[4]) ProductLog[-1,
  -(2 Log[4])/3 (1/(Sqrt[Pi] y))^(2/3)];

catData = {{2, 1, 1}, {1, 1, -1}, {1, 2, -1}};
catQ = Table[gammaLogCoefficient[catData, j], {j, 1, 6}];
Table[balancedInverseCoefficient[catQ, Log[4], -3/2, n],
  {n, 1, 3}]
Table[harmonicInverseCoefficient[m], {m, 1, 5}]
\end{lstlisting}

The intended first output is the three coefficients in \eqref{eq:catalan-inverse}; the second is \eqref{eq:harmonic-coeff}. Exact input constants should be kept exact until a final numerical evaluation. For very large targets, use high precision and work with logarithmic phases: constructing an enormous combinatorial number at machine precision and then taking its logarithm can lose the small corrections being studied.

A safe finite-order implementation has three separate precision controls: algebraic order within a sector, exponential action cutoff, and numerical working precision. They should not be represented by a single undocumented integer. For example, computing the involution amplitude through $x^{-5/2}$ does not resolve the $\ee^{-2\sqrt x}$ sector, and six total degrees in the Stirling-II inverse do not mean six successive dominance-ordered terms.

\section{Source and sequence guide}
\label{app:sources}

The following guide records the principal public definitions used in choosing the examples. OEIS entries are used for identification and exact sequence formulas, not as a substitute for the inverse proofs in the body of this article. The DLMF supplies standard analytic asymptotic inputs, and the explicit involution coefficients have the independent primary-paper check \cite{ekz}.

\begin{center}
\small
\begin{tabular}{P{51mm}P{35mm}P{56mm}}
\toprule
Object & Identifier or function & Role here\\\midrule
Catalan numbers & A000108; \texttt{CatalanNumber} & Balanced gamma interpolation\\
Ternary Fuss--Catalan & A001764 & Fixed-$p$ gamma family\\
Three-row rectangular tableaux & A005789 & Fixed-height rectangular specialization\\
Set partitions into $k$ blocks & \texttt{StirlingS2} & Exact finite exponential sum\\
Permutations with $k$ cycles & Signed \texttt{StirlingS1} & Unsigned gamma-coefficient interpolation\\
Motzkin numbers & A001006 & Two endpoint moments\\
Central Delannoy numbers & A001850 & Positive Legendre moment\\
Large Schr\"oder numbers & A006318 & Positive square-root moment\\
Involution numbers & A000085 & Two Gaussian saddles\\
Alternating permutations & A000111 & Parity-dependent Dirichlet sectors\\
All / connected labeled graphs & A006125 / A001187 & Formal component logarithm\\
Binary necklaces / Lyndon words & A000031 / A001037 & Divisor-action examples\\
Harmonic functions & \texttt{HarmonicNumber} & Slow-growth and endpoint inversion\\
Lambert branches & \texttt{ProductLog} & Exact core inversion\\\bottomrule
\end{tabular}
\end{center}

All sequence parameters and interpolation conventions used in the proofs are stated in their respective sections. The online references were consulted on September 4, 2026. URLs are included to make the reference trail reproducible; dynamically edited pages can subsequently change.

\begin{thebibliography}{99}
\small

\bibitem{base}
Vladimir Reshetnikov, \emph{Transseries: the polynomial--logarithmic calculus, series reversal at infinity, and the inversion of rapidly growing functions}, ProveIt repository, editable source and directory description.\newline
\url{https://github.com/VladimirReshetnikov/ProveIt/tree/main/Analysis/FabiusFunction/docs/semi-formalized-research-frontiers/drafts/series-and-transseries/Transseries_And_Inversion}

\bibitem{dlmfgamma}
NIST Digital Library of Mathematical Functions, \S5.11, \emph{Asymptotic Expansions} (gamma and digamma functions).\newline
\url{https://dlmf.nist.gov/5.11}

\bibitem{dlmfbernoulli}
NIST Digital Library of Mathematical Functions, \S24.2, \emph{Definitions and Generating Functions} (Bernoulli and Euler polynomials).\newline
\url{https://dlmf.nist.gov/24.2}

\bibitem{dlmflaplace}
NIST Digital Library of Mathematical Functions, \S2.3, \emph{Integrals of a Real Variable} (Watson's lemma and Laplace's method).\newline
\url{https://dlmf.nist.gov/2.3}

\bibitem{dlmfstirling}
NIST Digital Library of Mathematical Functions, \S26.8, \emph{Set Partitions: Stirling Numbers}.\newline
\url{https://dlmf.nist.gov/26.8}

\bibitem{dlmflegendre}
NIST Digital Library of Mathematical Functions, \S14.12, \emph{Integral Representations} (Legendre functions).\newline
\url{https://dlmf.nist.gov/14.12}

\bibitem{dlmfhurwitz}
NIST Digital Library of Mathematical Functions, \S25.11, \emph{Hurwitz Zeta Function}.\newline
\url{https://dlmf.nist.gov/25.11}

\bibitem{oeiscatalan}
The OEIS Foundation, \emph{A000108: Catalan numbers}.\newline
\url{https://oeis.org/A000108}

\bibitem{oeisfuss}
The OEIS Foundation, \emph{A001764: $\binom{3n}{n}/(2n+1)$, ternary trees}.\newline
\url{https://oeis.org/A001764}

\bibitem{oeistableaux}
The OEIS Foundation, \emph{A005789: standard Young tableaux of shape $(n,n,n)$}.\newline
\url{https://oeis.org/A005789}

\bibitem{oeismotzkin}
The OEIS Foundation, \emph{A001006: Motzkin numbers}.\newline
\url{https://oeis.org/A001006}

\bibitem{oeisdelannoy}
The OEIS Foundation, \emph{A001850: central Delannoy numbers}.\newline
\url{https://oeis.org/A001850}

\bibitem{oeisschroder}
The OEIS Foundation, \emph{A006318: large Schr\"oder numbers}.\newline
\url{https://oeis.org/A006318}

\bibitem{oeisinvolution}
The OEIS Foundation, \emph{A000085: involution numbers}.\newline
\url{https://oeis.org/A000085}

\bibitem{ekz}
Shalosh B. Ekhad, Manuel Kauers, and Doron Zeilberger,
\emph{The $O(1/n^{85})$ Asymptotic Expansion of OEIS Sequence A85}, 2024,
arXiv:2410.16334.\newline
\url{https://arxiv.org/abs/2410.16334}

\bibitem{oeiszigzag}
The OEIS Foundation, \emph{A000111: Euler or up/down numbers}.\newline
\url{https://oeis.org/A000111}

\bibitem{oeisallgraphs}
The OEIS Foundation, \emph{A006125: $2^{\binom n2}$, labeled graphs}.\newline
\url{https://oeis.org/A006125}

\bibitem{oeisconnected}
The OEIS Foundation, \emph{A001187: connected labeled graphs}.\newline
\url{https://oeis.org/A001187}

\bibitem{oeisnecklace}
The OEIS Foundation, \emph{A000031: binary necklaces}.\newline
\url{https://oeis.org/A000031}

\bibitem{oeislyndon}
The OEIS Foundation, \emph{A001037: primitive binary necklaces and Lyndon words}.\newline
\url{https://oeis.org/A001037}

\bibitem{wlcatalan}
Wolfram Research, \emph{CatalanNumber}, Wolfram Language Documentation.\newline
\url{https://reference.wolfram.com/language/ref/CatalanNumber.html}

\bibitem{wlstirling1}
Wolfram Research, \emph{StirlingS1}, Wolfram Language Documentation (signed first-kind convention).\newline
\url{https://reference.wolfram.com/language/ref/StirlingS1.html}

\bibitem{wlstirling2}
Wolfram Research, \emph{StirlingS2}, Wolfram Language Documentation.\newline
\url{https://reference.wolfram.com/language/ref/StirlingS2.html}

\bibitem{wlbeta}
Wolfram Research, \emph{DirichletBeta}, Wolfram Language Documentation.\newline
\url{https://reference.wolfram.com/language/ref/DirichletBeta.html}

\bibitem{wlharmonic}
Wolfram Research, \emph{HarmonicNumber}, Wolfram Language Documentation.\newline
\url{https://reference.wolfram.com/language/ref/HarmonicNumber.html}

\bibitem{wlproductlog}
Wolfram Research, \emph{ProductLog}, Wolfram Language Documentation.\newline
\url{https://reference.wolfram.com/language/ref/ProductLog.html}

\end{thebibliography}
\end{document}
